% Pass: Opus 5 (claude-opus-5), 2026-09-06 - first pass, unchecked against the
% sheets by a human. Its own context, no batch before it in this conversation.
%
% Book I, batch 8: sheets 85 to 96, leaves 71 to 81. Leaf 71 is the tail of
% batch 7 and leaf 75 is photographed twice, so the batch covers eleven leaves
% across twelve sheets.
%
% Six things worth knowing before reading it.
%
% The batch opens by paying batch 7's debt. Sheet 16955 is leaf 71 again with
% the Farm at Essex clipping turned back, and it shows the two blocks the
% clipping covered - Shore Acres and Barn at Essex. Only those two are
% transcribed here; the four entries batch 7 already gave at sheet 17836 are
% not repeated, and the two readings that complete a fifth are put in a note
% rather than into a second table.
%
% The subject changes three times, which batch 7's eight uniform leaves did
% not. Leaves 71 to 75 and 78 to 79 are the Cape Cod watercolours of 1929 to
% 1931 in the shape batch 7 established - a date, a price, a description, the
% amount net of Rehn's third, the date the cheque cleared. Leaves 76 and 77 cut
% into that run with two leaves of etchings at the Rehn Gallery, seventy-odd
% rows of $18 and $25 plates that belong with the etchings leaves at the front
% of the book. Then leaves 80 and 81 change register again: four oils at $1200
% to $5500, each with a large record drawing and two hundred words of
% description, and no ruled columns at all.
%
% Almost nothing is headed. Leaf 76 heads two columns, « Rec'd. » and « When
% rec'd. », and leaf 77 heads none; the watercolour leaves head none either.
% This is batch 5's régime rather than batch 6's and 7's, and no headings are
% supplied in \add{} anywhere in this batch.
%
% Leaf 75 is the happy hinged-clipping case a third time. On sheet 16910 a
% halftone of a house lies over the lower left of the leaf; on sheet 17138,
% which the Whitney calls « Page 75 - Verso », it has been turned back clear of
% the writing altogether. The transcription is at 17138, where the leaf is
% whole, and the clipping and a note stand at 16910.
%
% Leaf 81 is the unhappy case, and it costs the batch a block. The Camel's Hump
% halftone is pasted flat over the lower two thirds of the leaf and there is no
% second photograph of it anywhere in the Whitney's sequence. Six lines of
% pencil survive at its left edge, each one cut off at the clipping; they are
% given as far as they go and every \ill{} in them means « under paper ».
%
% Leaf 72's foot is not a work's record. Below two ordinary entries the leaf
% carries a running list headed « Unlisted Water Colors », where she tries to
% match pictures sold in 1955 to 1957 against titles she never entered, and
% twice asks the leaf itself whether two names are the same picture. It is
% transcribed as she set it down, in two columns, and the questions are hers.
\documentclass[11pt,a4paper]{article}
% The preamble shared by every transcription of the Hopper artist's ledgers.
%
% It rests on one idea, taken whole from the Grothendieck workbench this
% method comes from: **uncertainty compounds, it does not resolve.** A
% transcription that smooths over a doubtful figure has destroyed the only
% thing it was made for — a reader must be able to tell, at every point, what
% was on the leaf and what was supplied.
%
% A ledger sharpens that. In a manuscript of mathematics an invented word is a
% wrong word. Here an invented figure is a *sale that did not happen*, at a
% price nobody paid, to a buyer who never bought — and it will be cited,
% because a table looks like data in a way that prose does not.
%
% The same commands are recognised by `scripts/render.mjs`, which produces the
% site's reading view. A macro added here must be added there too, or rendering
% fails loudly, which is the wanted behaviour: a silently wrong rendering is
% worse than none.

% \ProvidesFile, not \ProvidesPackage: the transcriptions pull this in with
% % The preamble shared by every transcription of the Hopper artist's ledgers.
%
% It rests on one idea, taken whole from the Grothendieck workbench this
% method comes from: **uncertainty compounds, it does not resolve.** A
% transcription that smooths over a doubtful figure has destroyed the only
% thing it was made for — a reader must be able to tell, at every point, what
% was on the leaf and what was supplied.
%
% A ledger sharpens that. In a manuscript of mathematics an invented word is a
% wrong word. Here an invented figure is a *sale that did not happen*, at a
% price nobody paid, to a buyer who never bought — and it will be cited,
% because a table looks like data in a way that prose does not.
%
% The same commands are recognised by `scripts/render.mjs`, which produces the
% site's reading view. A macro added here must be added there too, or rendering
% fails loudly, which is the wanted behaviour: a silently wrong rendering is
% worse than none.

% \ProvidesFile, not \ProvidesPackage: the transcriptions pull this in with
% % The preamble shared by every transcription of the Hopper artist's ledgers.
%
% It rests on one idea, taken whole from the Grothendieck workbench this
% method comes from: **uncertainty compounds, it does not resolve.** A
% transcription that smooths over a doubtful figure has destroyed the only
% thing it was made for — a reader must be able to tell, at every point, what
% was on the leaf and what was supplied.
%
% A ledger sharpens that. In a manuscript of mathematics an invented word is a
% wrong word. Here an invented figure is a *sale that did not happen*, at a
% price nobody paid, to a buyer who never bought — and it will be cited,
% because a table looks like data in a way that prose does not.
%
% The same commands are recognised by `scripts/render.mjs`, which produces the
% site's reading view. A macro added here must be added there too, or rendering
% fails loudly, which is the wanted behaviour: a silently wrong rendering is
% worse than none.

% \ProvidesFile, not \ProvidesPackage: the transcriptions pull this in with
% \input{../preamble/hopper} rather than \usepackage, and \ProvidesPackage in
% an \input-ed file makes LaTeX warn on every compile that it "requested
% package `'" and got `hopper'. A warning nobody can act on is a warning
% everybody learns to scroll past, which is how the real ones get missed.
\ProvidesFile{hopper.sty}[2026/09/05 Transcriptions of the Hopper artist's ledgers]

\usepackage{array}
% longtable, not tabularx: a leaf of thirty-one exhibition lines must break
% across pages, and tabularx cannot be wrapped in an environment of our own —
% it scans the source for a literal \end{tabularx} and dies on \end{ledgertable}
% with « File ended while scanning use of \TX@get@body ». longtable has no such
% scan, and the wrapping column type below does the job tabularx's X would.
\usepackage{longtable}

% A wrapping column, `Y{0.3}` being three tenths of the measure.
%
% Prose columns must be declared this way and not as `l`. A table of `l`
% columns is set to its natural width, which can be wider than the page, and
% TeX issues **no warning at all** because nothing ever asked the row to fit:
% the first compile of Book I produced a page whose right-hand column ran off
% the paper with a clean log. A fixed-width column wraps, and if it still
% cannot fit, TeX says so — which is what `npm run pdf` then refuses to build
% over.
%
% Leave about 0.03 of the measure per column for the inter-column space: five
% columns can share 0.85, six can share 0.80. Getting it wrong is not a
% guessing game — `npm run pdf` reports the overfull box and how many points
% too wide it is.
\newcolumntype{Y}[1]{>{\raggedright\arraybackslash}p{#1\linewidth}}
\usepackage[english]{babel}
\usepackage{fontspec}
\usepackage{geometry}
\geometry{a4paper,margin=25mm,marginparwidth=28mm}
\usepackage{marginnote}
\usepackage{xcolor}
\usepackage{tikz}
% ulem, not soul. `soul`'s \st and \ul are built on a character-by-character
% scan that XeLaTeX's font machinery breaks: the document fails with an
% undefined control sequence rather than degrading. `normalem` stops it
% hijacking \emph, which the transcriptions use for underlining of Jo Hopper's
% own.
\usepackage[normalem]{ulem}
% ulem sets each underlined word in a box TeX can neither hyphenate nor break
% after, so a paragraph dense in \uncertain{} readings can reach a state where
% no legal breakpoint exists. \emergencystretch lets TeX loosen it on a third
% pass instead of overrunning: an overfull line in a transcription hides
% content.
\setlength{\emergencystretch}{3em}
\usepackage{hyperref}
\hypersetup{colorlinks=true,linkcolor=black,urlcolor=black,pdfborder={0 0 0}}

% --- Batch metadata ---------------------------------------------------------
% Taken as they stand by the rendering script, which shows them at the head of
% the reading view. They fix what the file claims to cover: without them a
% transcription floats unmoored in an archive of five hundred sheets.

\newcommand{\ledger}[1]{\gdef\hopLedger{#1}}
\newcommand{\ledgertitle}[1]{\gdef\hopLedgerTitle{#1}}
% A notebook of Josephine Hopper's, in place of a ledger: one file to a
% notebook, filed under transcripts/notebooks/, the argument the site's slug
% (« battle-of-wash-sq »). Sets the ledger to « notebooks » for everything
% that reads it, so the title block and the watermark behave as they do for a
% batch; the rendering checks the sheets against the notebook's own listing.
\newcommand{\notebook}[1]{\gdef\hopLedger{notebooks}\gdef\hopNotebook{#1}}
\gdef\hopNotebook{}
% The Whitney's accession number for the physical volume — 96.208 … 96.213.
% This is what a citation of the object cites.
\newcommand{\objectnumber}[1]{\gdef\hopObject{#1}}
% The batch this file covers. Leaving it out drops the « (batch N) » from the
% title block rather than printing a number that would be false.
\newcommand{\batch}[1]{\gdef\hopBatch{#1}}
\newcommand{\sheets}[2]{\gdef\hopFirst{#1}\gdef\hopLast{#2}}

% The Whitney's date range for the volume, copied verbatim from its resource
% record — never inferred here, never narrowed to the years this batch
% happens to cover.
\newcommand{\dating}[1]{\gdef\hopDating{#1}}

% The legal notice, printed in the title block and repeated as a diagonal
% watermark on every page of the PDF.
%
% This is not boilerplate and it is not optional. The ledgers are © Heirs of
% Josephine N. Hopper, licensed by the Artists Rights Society; the Whitney
% records the object rights as transferred to the Museum. A transcription
% reproduces Josephine Hopper's text. Every file must therefore say what it is
% on its own face, not only in a README that does not travel with the PDF.
% A file without this line gets no watermark, so omitting it is visible in the
% output rather than silent.
\newcommand{\watermark}[1]{\gdef\hopWatermark{#1}}

\gdef\hopLedger{?}\gdef\hopLedgerTitle{}\gdef\hopObject{}\gdef\hopBatch{}%
\gdef\hopFirst{?}\gdef\hopLast{?}\gdef\hopDating{}\gdef\hopWatermark{}

% --- Keywords ---------------------------------------------------------------
%
% Closing a transcription: the vocabulary under which to search for what these
% leaves record — dealers, exhibiting societies, media, places.
% `npm run manifest` extracts them to tag the whole ledger. This line is the
% single source of the tags; there is no tags file, so no tag can describe
% leaves nobody has read.
%
% It is the one thing in the file that is not on the paper, and it earns its
% place: a reader looking for Keppel needs some way in, and a tag written by
% whoever read the sheets is the only kind that cannot describe unread ones.
\newcommand{\keywords}[1]{\par\smallskip\noindent
  {\footnotesize\textsc{Keywords} --- \textit{#1}}\par}

% --- The anchor -------------------------------------------------------------

% The sheet begins here. Two arguments, and both are needed.
%
% #1 is the Whitney's **resource ref** — 16853 — which is the sheet's whole
% address: it is what the image URL is built from and what turns the facsimile
% as the transcript is scrolled. #2 is the number **written on the paper**, or
% empty where nobody wrote one, which is the case for every cover, flyleaf,
% index and loose insertion — about one sheet in nine.
%
% Keeping both is the whole of the numbering problem in this archive. Three
% sequences overlap: the Hoppers' own leaf numbers, the Whitney's order of
% photographs, and the resource refs, which are in no order at all. Only the
% first is citable and only the third is addressable, so a transcription that
% recorded one of them would be either uncitable or unopenable.
%
% `scripts/render.mjs` checks #2 against what `scripts/catalogue.mjs` parsed
% out of the Whitney's descriptor and fails on a disagreement, because a
% mistyped leaf number silently moves an entry to another year.
\newcommand{\sheet}[2]{%
  \marginnote{\footnotesize\color{gray!70}\textsf{\ifx\relax#2\relax---\else#2\fi}}%
  \label{sheet:#1}\ignorespaces}

% --- The critical apparatus -------------------------------------------------

% Illegible: never guessed. On a ledger leaf this is most often a figure, and
% a guessed figure is a fabricated transaction.
\newcommand{\ill}{\textcolor{red!60!black}{[\,\ldots\,]}}

% A doubtful reading: the word or figure is given, and flagged as uncertain.
%
% Underlining belongs to this macro and to nothing else. Jo Hopper underlines
% her own headings --- « Etchings », « Etching Plates », « Early oil » --- and
% those are set with \emph{} instead, because a reader must be able to tell a
% doubtful reading from her emphasis at a glance and the two cannot both be
% underlines. Which of them keeps the underline is decided by consequence: an
% emphasis shown as italic loses nothing, and a doubtful figure shown as a
% certain one loses everything.
\newcommand{\uncertain}[1]{\uline{#1}}

% An editorial addition — an expanded abbreviation, an implied dollar sign.
\newcommand{\add}[1]{\textcolor{blue!50!black}{[#1]}}

% Struck out in the book. Jo Hopper struck a great deal: a price revised, a
% buyer who withdrew, an exhibition that fell through. What she crossed out is
% part of the record and often the more interesting half of it.
\newcommand{\struck}[1]{\sout{#1}}

% The ink it is written in, where the ink says something.
%
% Book IV is a state machine and the colour is its status field: charges in
% black, receipts in red, the sums in pencil, a rule that holds from leaf 7 to
% the end of the volume. A transcription that records the words and the
% figures and nothing about their colour throws that field away, and the
% accounts then have to guess a receipt from the word « rec'd » --- which the
% volume stops writing on leaf 157. So the ink is recorded, per cell, as what
% is on the paper and not as what it means: \ink{red}{rec'd by check} says
% the words are red, and it is `scripts/accounts.mjs` that says red is money
% received. Black is the default and is not marked. The inks are the three
% the volume uses besides black; a fourth would be added here and in
% `scripts/render.mjs` and `scripts/tei.mjs` in the same commit.
% The file is \input, not \usepackage'd, so @ is not a letter here without saying so.
\makeatletter
\@namedef{hop@ink@red}{red!70!black}
\@namedef{hop@ink@pencil}{gray!55!black}
\@namedef{hop@ink@blue}{blue!55!black}
\newcommand{\ink}[2]{%
  \ifcsname hop@ink@#1\endcsname
    \textcolor{\csname hop@ink@#1\endcsname}{#2}%
  \else
    \PackageError{hopper}{\string\ink{#1}: unknown ink}{The inks are red, pencil and blue.}%
  \fi}
\makeatother

% A rule drawn across the money column above this row's figure.
%
% The writer rules off a run of items and writes their sum under the rule,
% and the rule is the whole of what says « this figure is a sum ». It stands
% at the head of the row the sum is on: \ruledoff & & 4355 & 00. On paper it
% is a line under the last two columns of the four a Book IV leaf has; if
% another volume ever rules a different column, the macro grows an argument.
\newcommand{\ruledoff}{\cline{3-4}}

% A diary entry begins. #1 is the date as she wrote it — « Mar. 12 »,
% « Sunday » — and #2 the date the transcriber assigns it, as a year, a month
% and a day (1947-03-12), or as much of one as the leaf allows. The first is
% what is on the paper; the second is the one editorial claim a diary needs
% that a ledger does not, because a ledger's leaf is its own address and an
% entry's date is its subject. Both are kept, apart, so the claim can be
% argued with. Used only in the notebooks; a ledger has no entries.
\newcommand{\entry}[2]{\par\medskip\noindent
  {\bfseries #1}\ifx\relax#2\relax\else\hspace{0.6em}{\footnotesize\color{gray!60!black}\textsf{#2}}\fi
  \par\nopagebreak\smallskip}

% The transcriber's note. Ours, not theirs.
\newcommand{\note}[1]{\par\smallskip{\small\color{gray!60!black}\textit{[#1]}}\par\smallskip}

% A marginal note **of theirs**, distinct from the body — and distinct from
% \note{} for exactly the reason \note{} exists.
\newcommand{\marginal}[1]{\par\smallskip\noindent\hspace{1em}%
  {\small\color{gray!60!black}#1}\par\smallskip}

% --- What is particular to this archive -------------------------------------

% Whose hand wrote it.
%
% This macro has no counterpart in the Grothendieck preamble, and it is the
% single most important thing in this one. These books were written by two
% people and annotated by more. Edward Hopper drew the work and wrote its
% title and size; Josephine Nivison Hopper wrote everything else — the
% descriptions, the anecdotes about people he refused to discuss, the prices,
% the buyers, the columns; and later hands, curatorial and unidentified, added
% to the leaves after 1967.
%
% A transcription that flattens the three has destroyed the archive's chief
% documentary value, which is that it records *two* working lives and the
% division of labour between them. #1 is `edward`, `jo`, `later` or
% `unidentified`; `unidentified` is the right answer far more often than a
% transcriber's confidence suggests, and it is not a failure to use it.
\newcommand{\hand}[2]{%
  \ifx\relax#1\relax#2\else
    \textcolor{gray!55!black}{\footnotesize\textsf{[#1]}}\,#2%
  \fi}

% Edward Hopper's ink record sketch stands here.
%
% The argument is **what is inscribed on or beside the sketch** — a title, a
% dimension, a note — and never a description of the image. Describing the
% drawing would be writing a new document: the sketch is on the site, one pane
% away, at the resolution the Whitney publishes, and prose about it competes
% with looking at it. Leave the argument empty where nothing is written.
% The mark is drawn rather than set from a symbol font: there is no
% mathematics anywhere in this archive, so no maths package is loaded, and
% $\square$ would be an undefined control sequence — which is exactly how it
% failed the first time the PDFs were compiled.
\newcommand{\sketchmark}{\raisebox{0.15ex}{\color{gray!50!black}%
  \tikz[baseline]{\draw[line width=0.4pt] (0,0) rectangle (1.1ex,1.1ex);}}}
\newcommand{\sketch}[1]{\par\smallskip\noindent
  {\small\color{gray!50!black}\sketchmark~\textsf{ink sketch}%
   \ifx\relax#1\relax\else\ --- #1\fi}\par\smallskip}

% Something printed and pasted to the leaf — a reproduction cut from a
% catalogue, a review, a dealer's label. The argument transcribes its printed
% text. These are the only places in the archive where the words are nobody's
% handwriting, and they date the leaf, which handwriting does not.
\newcommand{\clipping}[1]{\par\smallskip\noindent
  {\small\color{gray!50!black}\textsf{[clipping]}\ #1}\par\smallskip}

% A work's block opens here, under the title **exactly as the leaf gives it** —
% « Les Deux Pigeon » and « Les Poillus » stay misspelled, because the
% misspelling is Hopper's and is how the entry is found.
\newcommand{\work}[1]{\par\medskip\noindent{\bfseries #1}\par\nopagebreak\smallskip}

% --- The ruled table --------------------------------------------------------
%
% Jo Hopper ruled these columns herself and kept them for forty years; the
% table is not a presentation of the content, it *is* the content. So it is
% transcribed as a table and rendered as one, on screen and on paper, from one
% source.
%
% #1 is the column specification; #2 is the header row, `&`-separated like any
% other. A cell she left blank is left blank; a cell that cannot be read takes
% \ill{}.
%
% **Use `Y{...}` for any column that can hold a sentence** — an exhibition, a
% dealer's terms, a note — and `l`, `r` or `c` only for dates, figures and
% short names. See the note on \newcolumntype{Y} above for why that is not a
% matter of taste.
%
% The header row is emitted **bare**, between rules, and is deliberately not
% wrapped in \textbf{}. It cannot be: an `&` inside a braced group is not an
% alignment tab, so `\textbf{Date & Exhibitions}` fails at the first header
% with « Forbidden control sequence found while scanning use of
% \check@nocorr@ » — which is TeX saying, at some distance, that the tabs are
% in the wrong place. Rules carry the header instead, which is how a scholarly
% table sets one anyway. (The reading view has no such constraint and renders
% the same row as <th>.)
\newenvironment{ledgertable}[2]
  {\par\smallskip\small
   \begin{longtable}{#1}
   \hline
   #2 \\
   \hline
   \endhead}
  {\end{longtable}\par\smallskip}

% --- The appendix of notes --------------------------------------------------
%
% Everything above the appendix is on the paper. Nothing in it is.
%
% A note says what a named source says about a work these leaves record --- what
% the holding museum measures the canvas at, which other leaf carries the other
% half of a transaction --- and it is the only prose in the file that a reader
% cannot check against the photograph. On the site that is handled by putting
% the note beside a link to its source; on paper a link is not enough, because
% a PDF outlives the site it was downloaded from and travels away from it. So
% every claim here prints its source in full, name and address both, and every
% note prints the day it was written and what wrote it.
%
% The appendix is assembled by `scripts/pdf.mjs` from
% `src/content/work-notes.json` at compile time and is in no `.tex` file. That
% is the point of it. The transcription is the edition; a note is a later and
% weaker kind of writing *about* the edition, and keeping the two in one source
% file would eventually put a sentence somebody inferred inside the document
% whose whole claim on a reader is that nothing in it was inferred.
%
% Only the notes on works this batch's leaves name are printed, so the appendix
% follows the batch rather than repeating the archive, and a batch that has
% earned no notes yet gets no appendix at all --- which is a true statement
% about how much has been checked, and the same statement `work-notes.json`
% makes by being nearly empty.

% A URL is the one string in this archive that can run a hundred characters
% without a space in it, and an overfull box fails the build here because in a
% transcription an overfull box hides content. So the appendix tells TeX where
% a URL may break and sets its apparatus ragged right, rather than leaving the
% line to overrun in a file nobody recompiles.
\def\UrlBreaks{\do\/\do\-\do\.\do\_\do\?\do\&\do\=\do\+}
\Urlmuskip=0mu plus 1mu

\newcommand{\notesappendix}{\clearpage
  \par\noindent{\large\bfseries Notes on the works}\par\smallskip
  \noindent{\footnotesize\color{gray!60!black}\raggedright
    Not on the paper, and not part of the transcription. Each note below
    concerns a work named on a leaf of this batch and says what a named source
    says about it --- a museum's own catalogue record, or another leaf of these
    ledgers. Every sentence carries the source that states it, printed with its
    address so that it can be followed from a copy of this file alone. Where a
    ledger and a museum disagree the disagreement is printed rather than
    resolved.\par}\par\medskip}

% One work: the title as the leaves give it, then the note.
\newcommand{\worknote}[2]{\par\medskip\noindent{\bfseries #1}\par\nopagebreak
  \smallskip\noindent #2\par}

% One claim under it: what is asserted, and who asserts it.
\newcommand{\workclaim}[2]{\par\smallskip\noindent
  {\footnotesize\color{gray!60!black}\raggedright\hangindent=1.6em\hangafter=0
   \hspace{1em}#1\ --- #2\par}}

% Who wrote the note and when, closing the block. A note is a model's reading
% of retrieved records and is weaker evidence than anything above it; the file
% records that rather than letting the appendix pass for the edition.
\newcommand{\worknoteby}[2]{\par\smallskip\noindent
  {\scriptsize\color{gray!55!black}\hspace{1em}Written #1 by #2.\par}\par\smallskip}

% --- Title ------------------------------------------------------------------

% The watermark is drawn on the shipout background so it sits under the text of
% every page, diagonal and light enough to leave the record readable.
\AddToHook{shipout/background}{%
  \ifx\hopWatermark\empty\else
    \begin{tikzpicture}[remember picture, overlay]
      \node[rotate=52, scale=3.4, align=center, text=gray!18,
            font=\sffamily\bfseries]
        at (current page.center) {\hopWatermark};
    \end{tikzpicture}%
  \fi}

\AtBeginDocument{%
  \begin{center}
    {\large\bfseries Edward and Josephine Hopper --- \hopLedgerTitle}\\[2pt]
    \ifx\hopObject\empty\else
      {\small Whitney Museum of American Art, \hopObject}\\[4pt]
    \fi
    {\small sheets \hopFirst--\hopLast
      \ifx\hopBatch\empty\else\ (batch \hopBatch)\fi}\\[2pt]
    \ifx\hopDating\empty\else
      {\small\color{gray!60!black}The Whitney's dating: \hopDating}\\[2pt]
    \fi
    \ifx\hopWatermark\empty\else
      {\small\bfseries\color{red!45!gray}\hopWatermark}\\[2pt]
    \fi
    {\footnotesize\color{gray!60!black}Transcribed from the digitised sheets published by the
     Whitney Museum of American Art.\\
     \textcopyright{} Heirs of Josephine N. Hopper, licensed by Artists Rights Society (ARS),
     New York.\\
     Uncertain readings are underlined, illegible passages marked [\,\ldots\,],
     editorial additions bracketed.}
  \end{center}
  \vspace{1em}}

\endinput
 rather than \usepackage, and \ProvidesPackage in
% an \input-ed file makes LaTeX warn on every compile that it "requested
% package `'" and got `hopper'. A warning nobody can act on is a warning
% everybody learns to scroll past, which is how the real ones get missed.
\ProvidesFile{hopper.sty}[2026/09/05 Transcriptions of the Hopper artist's ledgers]

\usepackage{array}
% longtable, not tabularx: a leaf of thirty-one exhibition lines must break
% across pages, and tabularx cannot be wrapped in an environment of our own —
% it scans the source for a literal \end{tabularx} and dies on \end{ledgertable}
% with « File ended while scanning use of \TX@get@body ». longtable has no such
% scan, and the wrapping column type below does the job tabularx's X would.
\usepackage{longtable}

% A wrapping column, `Y{0.3}` being three tenths of the measure.
%
% Prose columns must be declared this way and not as `l`. A table of `l`
% columns is set to its natural width, which can be wider than the page, and
% TeX issues **no warning at all** because nothing ever asked the row to fit:
% the first compile of Book I produced a page whose right-hand column ran off
% the paper with a clean log. A fixed-width column wraps, and if it still
% cannot fit, TeX says so — which is what `npm run pdf` then refuses to build
% over.
%
% Leave about 0.03 of the measure per column for the inter-column space: five
% columns can share 0.85, six can share 0.80. Getting it wrong is not a
% guessing game — `npm run pdf` reports the overfull box and how many points
% too wide it is.
\newcolumntype{Y}[1]{>{\raggedright\arraybackslash}p{#1\linewidth}}
\usepackage[english]{babel}
\usepackage{fontspec}
\usepackage{geometry}
\geometry{a4paper,margin=25mm,marginparwidth=28mm}
\usepackage{marginnote}
\usepackage{xcolor}
\usepackage{tikz}
% ulem, not soul. `soul`'s \st and \ul are built on a character-by-character
% scan that XeLaTeX's font machinery breaks: the document fails with an
% undefined control sequence rather than degrading. `normalem` stops it
% hijacking \emph, which the transcriptions use for underlining of Jo Hopper's
% own.
\usepackage[normalem]{ulem}
% ulem sets each underlined word in a box TeX can neither hyphenate nor break
% after, so a paragraph dense in \uncertain{} readings can reach a state where
% no legal breakpoint exists. \emergencystretch lets TeX loosen it on a third
% pass instead of overrunning: an overfull line in a transcription hides
% content.
\setlength{\emergencystretch}{3em}
\usepackage{hyperref}
\hypersetup{colorlinks=true,linkcolor=black,urlcolor=black,pdfborder={0 0 0}}

% --- Batch metadata ---------------------------------------------------------
% Taken as they stand by the rendering script, which shows them at the head of
% the reading view. They fix what the file claims to cover: without them a
% transcription floats unmoored in an archive of five hundred sheets.

\newcommand{\ledger}[1]{\gdef\hopLedger{#1}}
\newcommand{\ledgertitle}[1]{\gdef\hopLedgerTitle{#1}}
% A notebook of Josephine Hopper's, in place of a ledger: one file to a
% notebook, filed under transcripts/notebooks/, the argument the site's slug
% (« battle-of-wash-sq »). Sets the ledger to « notebooks » for everything
% that reads it, so the title block and the watermark behave as they do for a
% batch; the rendering checks the sheets against the notebook's own listing.
\newcommand{\notebook}[1]{\gdef\hopLedger{notebooks}\gdef\hopNotebook{#1}}
\gdef\hopNotebook{}
% The Whitney's accession number for the physical volume — 96.208 … 96.213.
% This is what a citation of the object cites.
\newcommand{\objectnumber}[1]{\gdef\hopObject{#1}}
% The batch this file covers. Leaving it out drops the « (batch N) » from the
% title block rather than printing a number that would be false.
\newcommand{\batch}[1]{\gdef\hopBatch{#1}}
\newcommand{\sheets}[2]{\gdef\hopFirst{#1}\gdef\hopLast{#2}}

% The Whitney's date range for the volume, copied verbatim from its resource
% record — never inferred here, never narrowed to the years this batch
% happens to cover.
\newcommand{\dating}[1]{\gdef\hopDating{#1}}

% The legal notice, printed in the title block and repeated as a diagonal
% watermark on every page of the PDF.
%
% This is not boilerplate and it is not optional. The ledgers are © Heirs of
% Josephine N. Hopper, licensed by the Artists Rights Society; the Whitney
% records the object rights as transferred to the Museum. A transcription
% reproduces Josephine Hopper's text. Every file must therefore say what it is
% on its own face, not only in a README that does not travel with the PDF.
% A file without this line gets no watermark, so omitting it is visible in the
% output rather than silent.
\newcommand{\watermark}[1]{\gdef\hopWatermark{#1}}

\gdef\hopLedger{?}\gdef\hopLedgerTitle{}\gdef\hopObject{}\gdef\hopBatch{}%
\gdef\hopFirst{?}\gdef\hopLast{?}\gdef\hopDating{}\gdef\hopWatermark{}

% --- Keywords ---------------------------------------------------------------
%
% Closing a transcription: the vocabulary under which to search for what these
% leaves record — dealers, exhibiting societies, media, places.
% `npm run manifest` extracts them to tag the whole ledger. This line is the
% single source of the tags; there is no tags file, so no tag can describe
% leaves nobody has read.
%
% It is the one thing in the file that is not on the paper, and it earns its
% place: a reader looking for Keppel needs some way in, and a tag written by
% whoever read the sheets is the only kind that cannot describe unread ones.
\newcommand{\keywords}[1]{\par\smallskip\noindent
  {\footnotesize\textsc{Keywords} --- \textit{#1}}\par}

% --- The anchor -------------------------------------------------------------

% The sheet begins here. Two arguments, and both are needed.
%
% #1 is the Whitney's **resource ref** — 16853 — which is the sheet's whole
% address: it is what the image URL is built from and what turns the facsimile
% as the transcript is scrolled. #2 is the number **written on the paper**, or
% empty where nobody wrote one, which is the case for every cover, flyleaf,
% index and loose insertion — about one sheet in nine.
%
% Keeping both is the whole of the numbering problem in this archive. Three
% sequences overlap: the Hoppers' own leaf numbers, the Whitney's order of
% photographs, and the resource refs, which are in no order at all. Only the
% first is citable and only the third is addressable, so a transcription that
% recorded one of them would be either uncitable or unopenable.
%
% `scripts/render.mjs` checks #2 against what `scripts/catalogue.mjs` parsed
% out of the Whitney's descriptor and fails on a disagreement, because a
% mistyped leaf number silently moves an entry to another year.
\newcommand{\sheet}[2]{%
  \marginnote{\footnotesize\color{gray!70}\textsf{\ifx\relax#2\relax---\else#2\fi}}%
  \label{sheet:#1}\ignorespaces}

% --- The critical apparatus -------------------------------------------------

% Illegible: never guessed. On a ledger leaf this is most often a figure, and
% a guessed figure is a fabricated transaction.
\newcommand{\ill}{\textcolor{red!60!black}{[\,\ldots\,]}}

% A doubtful reading: the word or figure is given, and flagged as uncertain.
%
% Underlining belongs to this macro and to nothing else. Jo Hopper underlines
% her own headings --- « Etchings », « Etching Plates », « Early oil » --- and
% those are set with \emph{} instead, because a reader must be able to tell a
% doubtful reading from her emphasis at a glance and the two cannot both be
% underlines. Which of them keeps the underline is decided by consequence: an
% emphasis shown as italic loses nothing, and a doubtful figure shown as a
% certain one loses everything.
\newcommand{\uncertain}[1]{\uline{#1}}

% An editorial addition — an expanded abbreviation, an implied dollar sign.
\newcommand{\add}[1]{\textcolor{blue!50!black}{[#1]}}

% Struck out in the book. Jo Hopper struck a great deal: a price revised, a
% buyer who withdrew, an exhibition that fell through. What she crossed out is
% part of the record and often the more interesting half of it.
\newcommand{\struck}[1]{\sout{#1}}

% The ink it is written in, where the ink says something.
%
% Book IV is a state machine and the colour is its status field: charges in
% black, receipts in red, the sums in pencil, a rule that holds from leaf 7 to
% the end of the volume. A transcription that records the words and the
% figures and nothing about their colour throws that field away, and the
% accounts then have to guess a receipt from the word « rec'd » --- which the
% volume stops writing on leaf 157. So the ink is recorded, per cell, as what
% is on the paper and not as what it means: \ink{red}{rec'd by check} says
% the words are red, and it is `scripts/accounts.mjs` that says red is money
% received. Black is the default and is not marked. The inks are the three
% the volume uses besides black; a fourth would be added here and in
% `scripts/render.mjs` and `scripts/tei.mjs` in the same commit.
% The file is \input, not \usepackage'd, so @ is not a letter here without saying so.
\makeatletter
\@namedef{hop@ink@red}{red!70!black}
\@namedef{hop@ink@pencil}{gray!55!black}
\@namedef{hop@ink@blue}{blue!55!black}
\newcommand{\ink}[2]{%
  \ifcsname hop@ink@#1\endcsname
    \textcolor{\csname hop@ink@#1\endcsname}{#2}%
  \else
    \PackageError{hopper}{\string\ink{#1}: unknown ink}{The inks are red, pencil and blue.}%
  \fi}
\makeatother

% A rule drawn across the money column above this row's figure.
%
% The writer rules off a run of items and writes their sum under the rule,
% and the rule is the whole of what says « this figure is a sum ». It stands
% at the head of the row the sum is on: \ruledoff & & 4355 & 00. On paper it
% is a line under the last two columns of the four a Book IV leaf has; if
% another volume ever rules a different column, the macro grows an argument.
\newcommand{\ruledoff}{\cline{3-4}}

% A diary entry begins. #1 is the date as she wrote it — « Mar. 12 »,
% « Sunday » — and #2 the date the transcriber assigns it, as a year, a month
% and a day (1947-03-12), or as much of one as the leaf allows. The first is
% what is on the paper; the second is the one editorial claim a diary needs
% that a ledger does not, because a ledger's leaf is its own address and an
% entry's date is its subject. Both are kept, apart, so the claim can be
% argued with. Used only in the notebooks; a ledger has no entries.
\newcommand{\entry}[2]{\par\medskip\noindent
  {\bfseries #1}\ifx\relax#2\relax\else\hspace{0.6em}{\footnotesize\color{gray!60!black}\textsf{#2}}\fi
  \par\nopagebreak\smallskip}

% The transcriber's note. Ours, not theirs.
\newcommand{\note}[1]{\par\smallskip{\small\color{gray!60!black}\textit{[#1]}}\par\smallskip}

% A marginal note **of theirs**, distinct from the body — and distinct from
% \note{} for exactly the reason \note{} exists.
\newcommand{\marginal}[1]{\par\smallskip\noindent\hspace{1em}%
  {\small\color{gray!60!black}#1}\par\smallskip}

% --- What is particular to this archive -------------------------------------

% Whose hand wrote it.
%
% This macro has no counterpart in the Grothendieck preamble, and it is the
% single most important thing in this one. These books were written by two
% people and annotated by more. Edward Hopper drew the work and wrote its
% title and size; Josephine Nivison Hopper wrote everything else — the
% descriptions, the anecdotes about people he refused to discuss, the prices,
% the buyers, the columns; and later hands, curatorial and unidentified, added
% to the leaves after 1967.
%
% A transcription that flattens the three has destroyed the archive's chief
% documentary value, which is that it records *two* working lives and the
% division of labour between them. #1 is `edward`, `jo`, `later` or
% `unidentified`; `unidentified` is the right answer far more often than a
% transcriber's confidence suggests, and it is not a failure to use it.
\newcommand{\hand}[2]{%
  \ifx\relax#1\relax#2\else
    \textcolor{gray!55!black}{\footnotesize\textsf{[#1]}}\,#2%
  \fi}

% Edward Hopper's ink record sketch stands here.
%
% The argument is **what is inscribed on or beside the sketch** — a title, a
% dimension, a note — and never a description of the image. Describing the
% drawing would be writing a new document: the sketch is on the site, one pane
% away, at the resolution the Whitney publishes, and prose about it competes
% with looking at it. Leave the argument empty where nothing is written.
% The mark is drawn rather than set from a symbol font: there is no
% mathematics anywhere in this archive, so no maths package is loaded, and
% $\square$ would be an undefined control sequence — which is exactly how it
% failed the first time the PDFs were compiled.
\newcommand{\sketchmark}{\raisebox{0.15ex}{\color{gray!50!black}%
  \tikz[baseline]{\draw[line width=0.4pt] (0,0) rectangle (1.1ex,1.1ex);}}}
\newcommand{\sketch}[1]{\par\smallskip\noindent
  {\small\color{gray!50!black}\sketchmark~\textsf{ink sketch}%
   \ifx\relax#1\relax\else\ --- #1\fi}\par\smallskip}

% Something printed and pasted to the leaf — a reproduction cut from a
% catalogue, a review, a dealer's label. The argument transcribes its printed
% text. These are the only places in the archive where the words are nobody's
% handwriting, and they date the leaf, which handwriting does not.
\newcommand{\clipping}[1]{\par\smallskip\noindent
  {\small\color{gray!50!black}\textsf{[clipping]}\ #1}\par\smallskip}

% A work's block opens here, under the title **exactly as the leaf gives it** —
% « Les Deux Pigeon » and « Les Poillus » stay misspelled, because the
% misspelling is Hopper's and is how the entry is found.
\newcommand{\work}[1]{\par\medskip\noindent{\bfseries #1}\par\nopagebreak\smallskip}

% --- The ruled table --------------------------------------------------------
%
% Jo Hopper ruled these columns herself and kept them for forty years; the
% table is not a presentation of the content, it *is* the content. So it is
% transcribed as a table and rendered as one, on screen and on paper, from one
% source.
%
% #1 is the column specification; #2 is the header row, `&`-separated like any
% other. A cell she left blank is left blank; a cell that cannot be read takes
% \ill{}.
%
% **Use `Y{...}` for any column that can hold a sentence** — an exhibition, a
% dealer's terms, a note — and `l`, `r` or `c` only for dates, figures and
% short names. See the note on \newcolumntype{Y} above for why that is not a
% matter of taste.
%
% The header row is emitted **bare**, between rules, and is deliberately not
% wrapped in \textbf{}. It cannot be: an `&` inside a braced group is not an
% alignment tab, so `\textbf{Date & Exhibitions}` fails at the first header
% with « Forbidden control sequence found while scanning use of
% \check@nocorr@ » — which is TeX saying, at some distance, that the tabs are
% in the wrong place. Rules carry the header instead, which is how a scholarly
% table sets one anyway. (The reading view has no such constraint and renders
% the same row as <th>.)
\newenvironment{ledgertable}[2]
  {\par\smallskip\small
   \begin{longtable}{#1}
   \hline
   #2 \\
   \hline
   \endhead}
  {\end{longtable}\par\smallskip}

% --- The appendix of notes --------------------------------------------------
%
% Everything above the appendix is on the paper. Nothing in it is.
%
% A note says what a named source says about a work these leaves record --- what
% the holding museum measures the canvas at, which other leaf carries the other
% half of a transaction --- and it is the only prose in the file that a reader
% cannot check against the photograph. On the site that is handled by putting
% the note beside a link to its source; on paper a link is not enough, because
% a PDF outlives the site it was downloaded from and travels away from it. So
% every claim here prints its source in full, name and address both, and every
% note prints the day it was written and what wrote it.
%
% The appendix is assembled by `scripts/pdf.mjs` from
% `src/content/work-notes.json` at compile time and is in no `.tex` file. That
% is the point of it. The transcription is the edition; a note is a later and
% weaker kind of writing *about* the edition, and keeping the two in one source
% file would eventually put a sentence somebody inferred inside the document
% whose whole claim on a reader is that nothing in it was inferred.
%
% Only the notes on works this batch's leaves name are printed, so the appendix
% follows the batch rather than repeating the archive, and a batch that has
% earned no notes yet gets no appendix at all --- which is a true statement
% about how much has been checked, and the same statement `work-notes.json`
% makes by being nearly empty.

% A URL is the one string in this archive that can run a hundred characters
% without a space in it, and an overfull box fails the build here because in a
% transcription an overfull box hides content. So the appendix tells TeX where
% a URL may break and sets its apparatus ragged right, rather than leaving the
% line to overrun in a file nobody recompiles.
\def\UrlBreaks{\do\/\do\-\do\.\do\_\do\?\do\&\do\=\do\+}
\Urlmuskip=0mu plus 1mu

\newcommand{\notesappendix}{\clearpage
  \par\noindent{\large\bfseries Notes on the works}\par\smallskip
  \noindent{\footnotesize\color{gray!60!black}\raggedright
    Not on the paper, and not part of the transcription. Each note below
    concerns a work named on a leaf of this batch and says what a named source
    says about it --- a museum's own catalogue record, or another leaf of these
    ledgers. Every sentence carries the source that states it, printed with its
    address so that it can be followed from a copy of this file alone. Where a
    ledger and a museum disagree the disagreement is printed rather than
    resolved.\par}\par\medskip}

% One work: the title as the leaves give it, then the note.
\newcommand{\worknote}[2]{\par\medskip\noindent{\bfseries #1}\par\nopagebreak
  \smallskip\noindent #2\par}

% One claim under it: what is asserted, and who asserts it.
\newcommand{\workclaim}[2]{\par\smallskip\noindent
  {\footnotesize\color{gray!60!black}\raggedright\hangindent=1.6em\hangafter=0
   \hspace{1em}#1\ --- #2\par}}

% Who wrote the note and when, closing the block. A note is a model's reading
% of retrieved records and is weaker evidence than anything above it; the file
% records that rather than letting the appendix pass for the edition.
\newcommand{\worknoteby}[2]{\par\smallskip\noindent
  {\scriptsize\color{gray!55!black}\hspace{1em}Written #1 by #2.\par}\par\smallskip}

% --- Title ------------------------------------------------------------------

% The watermark is drawn on the shipout background so it sits under the text of
% every page, diagonal and light enough to leave the record readable.
\AddToHook{shipout/background}{%
  \ifx\hopWatermark\empty\else
    \begin{tikzpicture}[remember picture, overlay]
      \node[rotate=52, scale=3.4, align=center, text=gray!18,
            font=\sffamily\bfseries]
        at (current page.center) {\hopWatermark};
    \end{tikzpicture}%
  \fi}

\AtBeginDocument{%
  \begin{center}
    {\large\bfseries Edward and Josephine Hopper --- \hopLedgerTitle}\\[2pt]
    \ifx\hopObject\empty\else
      {\small Whitney Museum of American Art, \hopObject}\\[4pt]
    \fi
    {\small sheets \hopFirst--\hopLast
      \ifx\hopBatch\empty\else\ (batch \hopBatch)\fi}\\[2pt]
    \ifx\hopDating\empty\else
      {\small\color{gray!60!black}The Whitney's dating: \hopDating}\\[2pt]
    \fi
    \ifx\hopWatermark\empty\else
      {\small\bfseries\color{red!45!gray}\hopWatermark}\\[2pt]
    \fi
    {\footnotesize\color{gray!60!black}Transcribed from the digitised sheets published by the
     Whitney Museum of American Art.\\
     \textcopyright{} Heirs of Josephine N. Hopper, licensed by Artists Rights Society (ARS),
     New York.\\
     Uncertain readings are underlined, illegible passages marked [\,\ldots\,],
     editorial additions bracketed.}
  \end{center}
  \vspace{1em}}

\endinput
 rather than \usepackage, and \ProvidesPackage in
% an \input-ed file makes LaTeX warn on every compile that it "requested
% package `'" and got `hopper'. A warning nobody can act on is a warning
% everybody learns to scroll past, which is how the real ones get missed.
\ProvidesFile{hopper.sty}[2026/09/05 Transcriptions of the Hopper artist's ledgers]

\usepackage{array}
% longtable, not tabularx: a leaf of thirty-one exhibition lines must break
% across pages, and tabularx cannot be wrapped in an environment of our own —
% it scans the source for a literal \end{tabularx} and dies on \end{ledgertable}
% with « File ended while scanning use of \TX@get@body ». longtable has no such
% scan, and the wrapping column type below does the job tabularx's X would.
\usepackage{longtable}

% A wrapping column, `Y{0.3}` being three tenths of the measure.
%
% Prose columns must be declared this way and not as `l`. A table of `l`
% columns is set to its natural width, which can be wider than the page, and
% TeX issues **no warning at all** because nothing ever asked the row to fit:
% the first compile of Book I produced a page whose right-hand column ran off
% the paper with a clean log. A fixed-width column wraps, and if it still
% cannot fit, TeX says so — which is what `npm run pdf` then refuses to build
% over.
%
% Leave about 0.03 of the measure per column for the inter-column space: five
% columns can share 0.85, six can share 0.80. Getting it wrong is not a
% guessing game — `npm run pdf` reports the overfull box and how many points
% too wide it is.
\newcolumntype{Y}[1]{>{\raggedright\arraybackslash}p{#1\linewidth}}
\usepackage[english]{babel}
\usepackage{fontspec}
\usepackage{geometry}
\geometry{a4paper,margin=25mm,marginparwidth=28mm}
\usepackage{marginnote}
\usepackage{xcolor}
\usepackage{tikz}
% ulem, not soul. `soul`'s \st and \ul are built on a character-by-character
% scan that XeLaTeX's font machinery breaks: the document fails with an
% undefined control sequence rather than degrading. `normalem` stops it
% hijacking \emph, which the transcriptions use for underlining of Jo Hopper's
% own.
\usepackage[normalem]{ulem}
% ulem sets each underlined word in a box TeX can neither hyphenate nor break
% after, so a paragraph dense in \uncertain{} readings can reach a state where
% no legal breakpoint exists. \emergencystretch lets TeX loosen it on a third
% pass instead of overrunning: an overfull line in a transcription hides
% content.
\setlength{\emergencystretch}{3em}
\usepackage{hyperref}
\hypersetup{colorlinks=true,linkcolor=black,urlcolor=black,pdfborder={0 0 0}}

% --- Batch metadata ---------------------------------------------------------
% Taken as they stand by the rendering script, which shows them at the head of
% the reading view. They fix what the file claims to cover: without them a
% transcription floats unmoored in an archive of five hundred sheets.

\newcommand{\ledger}[1]{\gdef\hopLedger{#1}}
\newcommand{\ledgertitle}[1]{\gdef\hopLedgerTitle{#1}}
% A notebook of Josephine Hopper's, in place of a ledger: one file to a
% notebook, filed under transcripts/notebooks/, the argument the site's slug
% (« battle-of-wash-sq »). Sets the ledger to « notebooks » for everything
% that reads it, so the title block and the watermark behave as they do for a
% batch; the rendering checks the sheets against the notebook's own listing.
\newcommand{\notebook}[1]{\gdef\hopLedger{notebooks}\gdef\hopNotebook{#1}}
\gdef\hopNotebook{}
% The Whitney's accession number for the physical volume — 96.208 … 96.213.
% This is what a citation of the object cites.
\newcommand{\objectnumber}[1]{\gdef\hopObject{#1}}
% The batch this file covers. Leaving it out drops the « (batch N) » from the
% title block rather than printing a number that would be false.
\newcommand{\batch}[1]{\gdef\hopBatch{#1}}
\newcommand{\sheets}[2]{\gdef\hopFirst{#1}\gdef\hopLast{#2}}

% The Whitney's date range for the volume, copied verbatim from its resource
% record — never inferred here, never narrowed to the years this batch
% happens to cover.
\newcommand{\dating}[1]{\gdef\hopDating{#1}}

% The legal notice, printed in the title block and repeated as a diagonal
% watermark on every page of the PDF.
%
% This is not boilerplate and it is not optional. The ledgers are © Heirs of
% Josephine N. Hopper, licensed by the Artists Rights Society; the Whitney
% records the object rights as transferred to the Museum. A transcription
% reproduces Josephine Hopper's text. Every file must therefore say what it is
% on its own face, not only in a README that does not travel with the PDF.
% A file without this line gets no watermark, so omitting it is visible in the
% output rather than silent.
\newcommand{\watermark}[1]{\gdef\hopWatermark{#1}}

\gdef\hopLedger{?}\gdef\hopLedgerTitle{}\gdef\hopObject{}\gdef\hopBatch{}%
\gdef\hopFirst{?}\gdef\hopLast{?}\gdef\hopDating{}\gdef\hopWatermark{}

% --- Keywords ---------------------------------------------------------------
%
% Closing a transcription: the vocabulary under which to search for what these
% leaves record — dealers, exhibiting societies, media, places.
% `npm run manifest` extracts them to tag the whole ledger. This line is the
% single source of the tags; there is no tags file, so no tag can describe
% leaves nobody has read.
%
% It is the one thing in the file that is not on the paper, and it earns its
% place: a reader looking for Keppel needs some way in, and a tag written by
% whoever read the sheets is the only kind that cannot describe unread ones.
\newcommand{\keywords}[1]{\par\smallskip\noindent
  {\footnotesize\textsc{Keywords} --- \textit{#1}}\par}

% --- The anchor -------------------------------------------------------------

% The sheet begins here. Two arguments, and both are needed.
%
% #1 is the Whitney's **resource ref** — 16853 — which is the sheet's whole
% address: it is what the image URL is built from and what turns the facsimile
% as the transcript is scrolled. #2 is the number **written on the paper**, or
% empty where nobody wrote one, which is the case for every cover, flyleaf,
% index and loose insertion — about one sheet in nine.
%
% Keeping both is the whole of the numbering problem in this archive. Three
% sequences overlap: the Hoppers' own leaf numbers, the Whitney's order of
% photographs, and the resource refs, which are in no order at all. Only the
% first is citable and only the third is addressable, so a transcription that
% recorded one of them would be either uncitable or unopenable.
%
% `scripts/render.mjs` checks #2 against what `scripts/catalogue.mjs` parsed
% out of the Whitney's descriptor and fails on a disagreement, because a
% mistyped leaf number silently moves an entry to another year.
\newcommand{\sheet}[2]{%
  \marginnote{\footnotesize\color{gray!70}\textsf{\ifx\relax#2\relax---\else#2\fi}}%
  \label{sheet:#1}\ignorespaces}

% --- The critical apparatus -------------------------------------------------

% Illegible: never guessed. On a ledger leaf this is most often a figure, and
% a guessed figure is a fabricated transaction.
\newcommand{\ill}{\textcolor{red!60!black}{[\,\ldots\,]}}

% A doubtful reading: the word or figure is given, and flagged as uncertain.
%
% Underlining belongs to this macro and to nothing else. Jo Hopper underlines
% her own headings --- « Etchings », « Etching Plates », « Early oil » --- and
% those are set with \emph{} instead, because a reader must be able to tell a
% doubtful reading from her emphasis at a glance and the two cannot both be
% underlines. Which of them keeps the underline is decided by consequence: an
% emphasis shown as italic loses nothing, and a doubtful figure shown as a
% certain one loses everything.
\newcommand{\uncertain}[1]{\uline{#1}}

% An editorial addition — an expanded abbreviation, an implied dollar sign.
\newcommand{\add}[1]{\textcolor{blue!50!black}{[#1]}}

% Struck out in the book. Jo Hopper struck a great deal: a price revised, a
% buyer who withdrew, an exhibition that fell through. What she crossed out is
% part of the record and often the more interesting half of it.
\newcommand{\struck}[1]{\sout{#1}}

% The ink it is written in, where the ink says something.
%
% Book IV is a state machine and the colour is its status field: charges in
% black, receipts in red, the sums in pencil, a rule that holds from leaf 7 to
% the end of the volume. A transcription that records the words and the
% figures and nothing about their colour throws that field away, and the
% accounts then have to guess a receipt from the word « rec'd » --- which the
% volume stops writing on leaf 157. So the ink is recorded, per cell, as what
% is on the paper and not as what it means: \ink{red}{rec'd by check} says
% the words are red, and it is `scripts/accounts.mjs` that says red is money
% received. Black is the default and is not marked. The inks are the three
% the volume uses besides black; a fourth would be added here and in
% `scripts/render.mjs` and `scripts/tei.mjs` in the same commit.
% The file is \input, not \usepackage'd, so @ is not a letter here without saying so.
\makeatletter
\@namedef{hop@ink@red}{red!70!black}
\@namedef{hop@ink@pencil}{gray!55!black}
\@namedef{hop@ink@blue}{blue!55!black}
\newcommand{\ink}[2]{%
  \ifcsname hop@ink@#1\endcsname
    \textcolor{\csname hop@ink@#1\endcsname}{#2}%
  \else
    \PackageError{hopper}{\string\ink{#1}: unknown ink}{The inks are red, pencil and blue.}%
  \fi}
\makeatother

% A rule drawn across the money column above this row's figure.
%
% The writer rules off a run of items and writes their sum under the rule,
% and the rule is the whole of what says « this figure is a sum ». It stands
% at the head of the row the sum is on: \ruledoff & & 4355 & 00. On paper it
% is a line under the last two columns of the four a Book IV leaf has; if
% another volume ever rules a different column, the macro grows an argument.
\newcommand{\ruledoff}{\cline{3-4}}

% A diary entry begins. #1 is the date as she wrote it — « Mar. 12 »,
% « Sunday » — and #2 the date the transcriber assigns it, as a year, a month
% and a day (1947-03-12), or as much of one as the leaf allows. The first is
% what is on the paper; the second is the one editorial claim a diary needs
% that a ledger does not, because a ledger's leaf is its own address and an
% entry's date is its subject. Both are kept, apart, so the claim can be
% argued with. Used only in the notebooks; a ledger has no entries.
\newcommand{\entry}[2]{\par\medskip\noindent
  {\bfseries #1}\ifx\relax#2\relax\else\hspace{0.6em}{\footnotesize\color{gray!60!black}\textsf{#2}}\fi
  \par\nopagebreak\smallskip}

% The transcriber's note. Ours, not theirs.
\newcommand{\note}[1]{\par\smallskip{\small\color{gray!60!black}\textit{[#1]}}\par\smallskip}

% A marginal note **of theirs**, distinct from the body — and distinct from
% \note{} for exactly the reason \note{} exists.
\newcommand{\marginal}[1]{\par\smallskip\noindent\hspace{1em}%
  {\small\color{gray!60!black}#1}\par\smallskip}

% --- What is particular to this archive -------------------------------------

% Whose hand wrote it.
%
% This macro has no counterpart in the Grothendieck preamble, and it is the
% single most important thing in this one. These books were written by two
% people and annotated by more. Edward Hopper drew the work and wrote its
% title and size; Josephine Nivison Hopper wrote everything else — the
% descriptions, the anecdotes about people he refused to discuss, the prices,
% the buyers, the columns; and later hands, curatorial and unidentified, added
% to the leaves after 1967.
%
% A transcription that flattens the three has destroyed the archive's chief
% documentary value, which is that it records *two* working lives and the
% division of labour between them. #1 is `edward`, `jo`, `later` or
% `unidentified`; `unidentified` is the right answer far more often than a
% transcriber's confidence suggests, and it is not a failure to use it.
\newcommand{\hand}[2]{%
  \ifx\relax#1\relax#2\else
    \textcolor{gray!55!black}{\footnotesize\textsf{[#1]}}\,#2%
  \fi}

% Edward Hopper's ink record sketch stands here.
%
% The argument is **what is inscribed on or beside the sketch** — a title, a
% dimension, a note — and never a description of the image. Describing the
% drawing would be writing a new document: the sketch is on the site, one pane
% away, at the resolution the Whitney publishes, and prose about it competes
% with looking at it. Leave the argument empty where nothing is written.
% The mark is drawn rather than set from a symbol font: there is no
% mathematics anywhere in this archive, so no maths package is loaded, and
% $\square$ would be an undefined control sequence — which is exactly how it
% failed the first time the PDFs were compiled.
\newcommand{\sketchmark}{\raisebox{0.15ex}{\color{gray!50!black}%
  \tikz[baseline]{\draw[line width=0.4pt] (0,0) rectangle (1.1ex,1.1ex);}}}
\newcommand{\sketch}[1]{\par\smallskip\noindent
  {\small\color{gray!50!black}\sketchmark~\textsf{ink sketch}%
   \ifx\relax#1\relax\else\ --- #1\fi}\par\smallskip}

% Something printed and pasted to the leaf — a reproduction cut from a
% catalogue, a review, a dealer's label. The argument transcribes its printed
% text. These are the only places in the archive where the words are nobody's
% handwriting, and they date the leaf, which handwriting does not.
\newcommand{\clipping}[1]{\par\smallskip\noindent
  {\small\color{gray!50!black}\textsf{[clipping]}\ #1}\par\smallskip}

% A work's block opens here, under the title **exactly as the leaf gives it** —
% « Les Deux Pigeon » and « Les Poillus » stay misspelled, because the
% misspelling is Hopper's and is how the entry is found.
\newcommand{\work}[1]{\par\medskip\noindent{\bfseries #1}\par\nopagebreak\smallskip}

% --- The ruled table --------------------------------------------------------
%
% Jo Hopper ruled these columns herself and kept them for forty years; the
% table is not a presentation of the content, it *is* the content. So it is
% transcribed as a table and rendered as one, on screen and on paper, from one
% source.
%
% #1 is the column specification; #2 is the header row, `&`-separated like any
% other. A cell she left blank is left blank; a cell that cannot be read takes
% \ill{}.
%
% **Use `Y{...}` for any column that can hold a sentence** — an exhibition, a
% dealer's terms, a note — and `l`, `r` or `c` only for dates, figures and
% short names. See the note on \newcolumntype{Y} above for why that is not a
% matter of taste.
%
% The header row is emitted **bare**, between rules, and is deliberately not
% wrapped in \textbf{}. It cannot be: an `&` inside a braced group is not an
% alignment tab, so `\textbf{Date & Exhibitions}` fails at the first header
% with « Forbidden control sequence found while scanning use of
% \check@nocorr@ » — which is TeX saying, at some distance, that the tabs are
% in the wrong place. Rules carry the header instead, which is how a scholarly
% table sets one anyway. (The reading view has no such constraint and renders
% the same row as <th>.)
\newenvironment{ledgertable}[2]
  {\par\smallskip\small
   \begin{longtable}{#1}
   \hline
   #2 \\
   \hline
   \endhead}
  {\end{longtable}\par\smallskip}

% --- The appendix of notes --------------------------------------------------
%
% Everything above the appendix is on the paper. Nothing in it is.
%
% A note says what a named source says about a work these leaves record --- what
% the holding museum measures the canvas at, which other leaf carries the other
% half of a transaction --- and it is the only prose in the file that a reader
% cannot check against the photograph. On the site that is handled by putting
% the note beside a link to its source; on paper a link is not enough, because
% a PDF outlives the site it was downloaded from and travels away from it. So
% every claim here prints its source in full, name and address both, and every
% note prints the day it was written and what wrote it.
%
% The appendix is assembled by `scripts/pdf.mjs` from
% `src/content/work-notes.json` at compile time and is in no `.tex` file. That
% is the point of it. The transcription is the edition; a note is a later and
% weaker kind of writing *about* the edition, and keeping the two in one source
% file would eventually put a sentence somebody inferred inside the document
% whose whole claim on a reader is that nothing in it was inferred.
%
% Only the notes on works this batch's leaves name are printed, so the appendix
% follows the batch rather than repeating the archive, and a batch that has
% earned no notes yet gets no appendix at all --- which is a true statement
% about how much has been checked, and the same statement `work-notes.json`
% makes by being nearly empty.

% A URL is the one string in this archive that can run a hundred characters
% without a space in it, and an overfull box fails the build here because in a
% transcription an overfull box hides content. So the appendix tells TeX where
% a URL may break and sets its apparatus ragged right, rather than leaving the
% line to overrun in a file nobody recompiles.
\def\UrlBreaks{\do\/\do\-\do\.\do\_\do\?\do\&\do\=\do\+}
\Urlmuskip=0mu plus 1mu

\newcommand{\notesappendix}{\clearpage
  \par\noindent{\large\bfseries Notes on the works}\par\smallskip
  \noindent{\footnotesize\color{gray!60!black}\raggedright
    Not on the paper, and not part of the transcription. Each note below
    concerns a work named on a leaf of this batch and says what a named source
    says about it --- a museum's own catalogue record, or another leaf of these
    ledgers. Every sentence carries the source that states it, printed with its
    address so that it can be followed from a copy of this file alone. Where a
    ledger and a museum disagree the disagreement is printed rather than
    resolved.\par}\par\medskip}

% One work: the title as the leaves give it, then the note.
\newcommand{\worknote}[2]{\par\medskip\noindent{\bfseries #1}\par\nopagebreak
  \smallskip\noindent #2\par}

% One claim under it: what is asserted, and who asserts it.
\newcommand{\workclaim}[2]{\par\smallskip\noindent
  {\footnotesize\color{gray!60!black}\raggedright\hangindent=1.6em\hangafter=0
   \hspace{1em}#1\ --- #2\par}}

% Who wrote the note and when, closing the block. A note is a model's reading
% of retrieved records and is weaker evidence than anything above it; the file
% records that rather than letting the appendix pass for the edition.
\newcommand{\worknoteby}[2]{\par\smallskip\noindent
  {\scriptsize\color{gray!55!black}\hspace{1em}Written #1 by #2.\par}\par\smallskip}

% --- Title ------------------------------------------------------------------

% The watermark is drawn on the shipout background so it sits under the text of
% every page, diagonal and light enough to leave the record readable.
\AddToHook{shipout/background}{%
  \ifx\hopWatermark\empty\else
    \begin{tikzpicture}[remember picture, overlay]
      \node[rotate=52, scale=3.4, align=center, text=gray!18,
            font=\sffamily\bfseries]
        at (current page.center) {\hopWatermark};
    \end{tikzpicture}%
  \fi}

\AtBeginDocument{%
  \begin{center}
    {\large\bfseries Edward and Josephine Hopper --- \hopLedgerTitle}\\[2pt]
    \ifx\hopObject\empty\else
      {\small Whitney Museum of American Art, \hopObject}\\[4pt]
    \fi
    {\small sheets \hopFirst--\hopLast
      \ifx\hopBatch\empty\else\ (batch \hopBatch)\fi}\\[2pt]
    \ifx\hopDating\empty\else
      {\small\color{gray!60!black}The Whitney's dating: \hopDating}\\[2pt]
    \fi
    \ifx\hopWatermark\empty\else
      {\small\bfseries\color{red!45!gray}\hopWatermark}\\[2pt]
    \fi
    {\footnotesize\color{gray!60!black}Transcribed from the digitised sheets published by the
     Whitney Museum of American Art.\\
     \textcopyright{} Heirs of Josephine N. Hopper, licensed by Artists Rights Society (ARS),
     New York.\\
     Uncertain readings are underlined, illegible passages marked [\,\ldots\,],
     editorial additions bracketed.}
  \end{center}
  \vspace{1em}}

\endinput


\ledger{book-i}
\ledgertitle{Artist's ledger --- Book I}
\objectnumber{96.208}
\batch{8}
\sheets{85}{96}
\dating{1913--1963}
\watermark{Demonstration edition}

\begin{document}

\section{Leaf 71 again: the blocks under the Farm at Essex clipping}

\sheet{16955}{71}

\note{this is the recto of leaf 71 photographed a second time, with the Farm at
Essex halftone turned back off the leaf. The Whitney calls it « Page 71 -
Verso »; what it shows is the recto again, whole. Batch 7 transcribed this leaf
at sheet 17836, where the clipping lay flat over the lower left and covered two
entire blocks. Those two blocks are given below and nothing else is: the four
entries already transcribed at 17836 are not repeated here.}

\work{Shore Acres}

\begin{ledgertable}{Y{0.10}Y{0.06}Y{0.44}Y{0.09}Y{0.14}}{ & & & & }
 & 500 & Shore Acres 16 x 25 \quad (Concrete Road) listed name. \quad Dense green trees, a few tree trunks defined in the mass. Fields to R. \quad \struck{\ill{}} \quad \struck{Mr. +} Mrs. Pack - Arthur \struck{Pack}. of Princeton, N.J. \quad \emph{Returned} toward purchase of Camel's Hump, big oil - fall 1931. \quad \uncertain{Mrs. Arthur Fleishman} . 750 - 1/3 = 500. & 333 1/3 & Nov. 10, 30 \quad 8'' \quad Jan. 11, 1958. \\
 & 500 & Barn at Essex 16 x 25 \quad Beauty \uncertain{one}. \quad Big white barn, green lawn in front - palish. palish sky, grey day (white day) \quad like atmosphere moist, but not fog. \quad Blue green grass in front \quad 750 - 1/3 = 500. & & Dec. 29, 51. \\
\end{ledgertable}

\marginal{\ill{}oncrete \\ Road '' \quad --- in the left margin beside the Shore
Acres block, the first letter cut away by the edge of the leaf}

\note{the struck word before « Mrs. Pack » is in pencil and heavily crossed;
« Painted » will support it and is not offered as a reading. « Mr. + » at the
head of that line and the second « Pack » are struck in ink. The buyer of
January 1958 is written in a blunt pencil that will not resolve: « Fleishman »
is the spelling leaf 58 gives for what may be the same family, and is offered
here rather than read. The word after « Beauty » in the Barn at Essex block is
finished by an ink blot.}

\note{two readings that complete an entry batch 7 gave as \ill{} at sheet
17836, and are set down here rather than transcribed a second time. Under
« Cleveland Museum - Mar. 1930. » in the Light House Village block the leaf
carries, in pencil, « for 400 \uncertain{vs.} 500 », and below it in ink
« Cleveland Museum \quad 400 - 1/3 ». The « 333 » of the Shore Acres row is
written over an earlier figure, as batch 7 noted, and the « 8'' » after the
date is as written; its sense is still not established.}

\section{Leaf 72: Pemaquid Light, House with a Vine, and the unlisted water colors}

\sheet{18117}{72}

\begin{ledgertable}{Y{0.10}Y{0.06}Y{0.44}Y{0.09}Y{0.14}}{ & & & & }
Dec. 12, 29 & 400 & Pemaquid Light \quad 14 x 20. \quad Pale color sky + light house - group to L. - keepers + friends. \quad Mr. Wheelright \quad Dec. 19. - Has Rock and Harbor - \quad Robt. Wheelwright \quad 400 - 1/3 \quad Chestnut Hill, Philadelphia. & 266 2/3 & Mar. 5, 31. \\
Dec. 12, 29 & 400 & House with a Vine \quad 14 x 20. \quad A beauty. Heavy vine over door of house. \quad Cape Elizabeth . E. H. favorite.. \quad Dec. 1945. \quad Mrs. Brewster Jennings \quad Mrs. Brewster Jennings - Ap. 4th, 46 . 400 - 1/3 = 266 2/3. & & Ap. 4, 1946. \\
\end{ledgertable}

\marginal{Yankee of Philadelphia \quad --- above « Dec. 19. » in the Pemaquid
Light block}

\marginal{? \quad --- above « Brewster » in the ink line of the House with a
Vine block}

\note{« Wheelright » in the pencil line and « Wheelwright » in the ink line
below it are two spellings of one name, four lines apart. « Dec. 19. - Has Rock
and Harbor - » is in pencil and was written before the ink line that repeats
the buyer.}

\hand{jo}{These all paid for - 1955 + 1957.}

\hand{jo}{\emph{Unlisted Water Colors} - }

\begin{ledgertable}{Y{0.42}Y{0.42}}{ & }
Court Hill \quad 1930 or thereabouts. & Sold to Ferdinand Lane, N.Y.C. - Xmas 1955. \\
 & 600 - 1/3 - Then returned for House on Hill Top - Ferdinand Davis \quad 400. \quad from \uncertain{Gally} Feb. 14, 1957. \quad 600 - 1/3 = \\
House on Beach, Gloucester - (not listed) & Bgt. by Fraad - Allied N.Y. Service ? - 550 - 1/3 = 366 2/3 \quad Feb. 14, 57. \\
Ryder Beach Road - & \ill{} Edwards . 600 - 1/3 = 400 \quad July 8, '57. \\
Harbor Shore, Rockland - & Burton . \quad '' \quad '' \quad '' \quad '' \quad '' \quad '' \\
Could this be Shipyard, Rockland - schooner Butterfly? \quad p. 64. & \\
Dune at Truro * & Stein - 900 - 1/3 = 600 \quad July 8, '57. \\
Could this be Dune with Green Top p. 74 - Curious object on top of hill ? . p. 74. & \\
Court Hill later sold to & Mlodinich \quad 900 - 1/3 = 600 \quad July 8, 57. \\
The Captain's House - Gloucester, next to Our Lady of Good Voyage - 1924, '26 \quad 14 X 20. & \\
Matted. & \\
Pretty Penny = drawing, matted. Preliminary to painting of Helen Hayes House, Nyack, N.Y. \quad 1939. Book II & \\
Owned by Jo Hopper? & \\
\end{ledgertable}

\marginal{* NO \quad Listed as sold. \quad --- ringed, at the right of the Dune
at Truro row, with a brace drawn to it}

\note{this block heads no column and rules none. What is set as two columns
here stands on the leaf as a run of underlined titles at the left with the
buyer, the terms and the date beginning at a consistent point to their right;
the two « Could this be » lines and the last four lines run the full width and
take the left cell alone. « Court Hill » is written as one word both times.
Several rows of this block carry pencil that has been erased and is not
transcribed. « Ferdinand Lane » and « Ferdinand Davis » are three lines apart
and both stand. The word before « Edwards » is erased pencil under the ink.}

\section{Leaf 73: Beginning of Cape Cod pictures, 1930}

\sheet{18145}{73}

\hand{jo}{Beginning of Cape Cod pictures 1930.}

\hand{jo}{Water colors Summer 1930 - South Truro. Cape Cod.}

\begin{ledgertable}{Y{0.10}Y{0.06}Y{0.44}Y{0.09}Y{0.14}}{ & & & & }
Oct. 10, 30 & 400 & North Truro Station \quad 14 x 20 \quad Car tracks across front. Name of station smeared \quad Mr. Howard. Etna Paper Co. Dayton O. Dec. 1930. \quad 400 - 1/3 & 266 2/3 & Jan. 19, 32. \\
'' & 400 & South Truro Post Office I \quad 14 X 20. \quad Little red P.O., trim little round tree at side \quad Our house sticking up behind trees \quad Mrs. Ruth Maguire \quad 400 - 1/3 . & 266 2/3 & May 27, 57. \\
'' & 400 & Truro Station Coal Box \quad 14 x 20 \quad Rather bare. \quad Wadsworth Atheneum \quad 400 - 1/3 \quad Hartford, Conn. & 266.67 & Dec. 31, 34. \\
'' & 400 & S. Truro P.O. II \quad 14 x 20 \quad Red P.O.; full blown. dense looking tree + end of Cobb house \quad This transfer accredited to House at Eastham. Vol. II \quad ap. 7, 50. \quad Alfred L. Hydeman of York, Pa., vice pres. P. Wiest Dept. Store. \quad 400 - 1/3 & 266.67 & July 22 \quad 1950. \\
'' & 400 & Burly Cobb's House \quad 14 x 20 . \quad Rear view on tops of roofs in light - house dark red - planes of roofs the interesting thing. \quad Ed. Root. Oct. \quad Edward W. Root \quad 400 - 1/3 & 266 2/3 & Mar. 5, 31. \\
'' & 400 & Burly Cobb Hen Coop + barn \quad 14 x 20 \quad Roof of the hen coop - shingles make the picture; hills beyond \quad Mrs. McGuire \quad Jan. 41. \quad Mrs. McGuire \quad 400 - 1/3 = 266 2/3 & & Ap. 19, 41. \\
\end{ledgertable}

\sketch{THE COAL BOX}

\sketch{THE HEN COOP}

\marginal{Visit with Edward + Grace Root at Hamilton College before starting E.
to Cape Cod. \quad --- in pencil, between the two heading lines}

\marginal{Roof of \quad --- above « Our house sticking up behind trees », with a
caret}

\marginal{1937 \quad --- in pencil above « May 27, 57 »}

\marginal{(died) \quad --- in pencil below « Mrs. Ruth Maguire »}

\marginal{Returned \uncertain{(in us.)} Ap. 7, 1956 as part payment for House at
Eastham \quad 750 - 400 Truro P.O.}

\marginal{The S. Truro P.O. sold to Geo. Greenspan - Ap. 7, 56 for that same
400 . which goes to E. H. as part payment of House at Eastham}

\note{the two marginal lines about the Eastham transfer are written in pencil
straight across the leaf below the S. Truro P.O. II block and belong to it. Both
sketches are drawn in a ruled box beside the entry and each carries its title in
drawn pencil capitals; the same words stand again, faintly, in the margin at the
left of each. « Maguire » in the second row and « McGuire » in the last are two
spellings of one name on one leaf. « Atheneum » here is a third form of the
Wadsworth, against « Antheneum » and « Athenaeum » on leaves 64 and 66. The
Howard line was written first in pencil and then gone over in ink, with « etna »
still standing in pencil above it. « Ruth » and « Ed. Root. Oct. » are pencil
under the ink of their rows.}

\section{Leaf 74: Watercolors Summer 1930, South Truro, continued}

\sheet{17216}{74}

\begin{ledgertable}{Y{0.10}Y{0.06}Y{0.44}Y{0.09}Y{0.14}}{ & & & & }
Oct. 10, 30 & 400 & Wellfleet Bridge \quad 14 x 20 \quad Bridge effect. bright blue waters; foreground dark green tallish grass, bright yellow light across top of it. \quad Mr. Huntington, of Hartford \quad Oct. 15 \quad Robt. W. Huntington of Hartford. & 266 2/3 & Mar. 5, 31. \\
'' & 500 & Highland Light. N. Truro \quad 16 X 25 \quad Light house + structures belonging to it. Explosive sky - white in back of building. ragged cut of white, blue beyond. Sultry colored ground + road. \quad Fogg Mus. Jan. 1931. & 333 1/3 & Jan. 19, 32. \\
'' & 500 & Rich's House \quad 16 X 25 \quad Big square white house with sturdy little tree at R. side. fence + field. Wind mill at R. \quad For David Rockefeller - thru Weyhe Gal. \quad 750 - 1/3 = 500. & & June 16, 53. \\
'' & 500 & Dune with Green top \quad 16 x 25 \quad Yellow sand, green growth over top. Funny house thing on top. Very blue sky with strip of white cloud. \quad Stein - 900 - 1/3 = 600 \quad Paid July 8, '57 \quad if this the right picture \quad Mrs. Bliss Parkinson \quad 1000 - 1/3 . & & Dec 30, 1958 \\
'' & & Geo. \uncertain{Tennure} (spelling?) Lenox - for 400 thru Mrs. Albert Sterner's Art Room at Lenox \quad obtained from Rehn Gallery \quad 400 - 1/3 - & 266 2/3 & Jan. 8, 34. \\
'' & 750 & House in Provincetown \quad 20 X 25 \quad White house with 3 gables, front in strong shadow. Balcony along 2'' story front. Little pavellion in yard at side, front of picture. Hedge along \uncertain{front} of house. \quad Amer. Federation State Dept. for Exhibition to China \quad 450 - 1/3 = 300. & & Feb. 3rd, 1947. \\
\end{ledgertable}

\marginal{Dune at Truro? \quad --- in pencil above « Dune with Green top »}

\marginal{NO \quad --- underlined, in the money column of the Dune with Green
top row, above « Dec 30, 1958 »}

\note{the fifth row has no title: it is a second sale line for the picture
above it or for none, and stands in the leaf exactly where it is given here.
« (spelling?) » is hers, written after the buyer's name, so the name is
doubtful on the leaf and not only to a reader. The last word of the House in
Provincetown description is written over an earlier one; « front » is offered
rather than read. « pavellion » is as written. A flap of paper covers the
fore-edge of this photograph beside the Dune with Green top row, and « 1958 »
in pencil below the ink date is the only thing legible under it.}

\section{Leaf 75 with the clipping down}

\sheet{16910}{75}

\clipping{a halftone of a watercolour --- a light house with a picket fence, a
telephone pole at its left, low hills behind --- printed on a cream leaf and
captioned « Edward Hopper » in italic below, hinged at its left edge and lying
over the lower left of the leaf.}

\note{the clipping covers the Whitney exchange lines, the whole of the Dead
Tree + Side of Lombard House block and the left of the Roofs of the Cobb Barn
block. The Whitney photographed this leaf a second time with the clipping
turned back --- sheet 17138, which it calls « Page 75 - Verso » and which is
the recto again --- and there the leaf is whole. The transcription is given at
17138. Nothing is transcribed here, so that no figure is counted twice.}

\section{Leaf 75: Methodist Church Tower, and Summer 1931 South Truro}

\sheet{17138}{75}

\begin{ledgertable}{Y{0.10}Y{0.06}Y{0.44}Y{0.09}Y{0.14}}{ & & & & }
Oct. 10, '30. & 750 & Methodist Church Tower \quad 20 x 25. \quad Mottled, windy swirl of torn clouds in back of tower. In front house with rusty red + yellow tiled roof. Foremost house very heavy moulding under rafters + table to sit on strongly painted. Rather sophisticated looking tower. Tower greyish in shadow. \quad Provincetown. \quad Sheaffer - Oct. 750 \quad Leslie G. Sheafer \quad 750 - 1/3 \quad \uncertain{Seoley} Green St. \quad 1951 ± Sold to Wadsworth Atheneum, Hartford Conn. & 500. & Mar. 5, 31. \\
\end{ledgertable}

\hand{jo}{\emph{Summer 1931 South Truro. + Truro.}}

\begin{ledgertable}{Y{0.10}Y{0.06}Y{0.44}Y{0.09}Y{0.14}}{ & & & & }
Oct. 1931. & & Rich's Barn \quad 20 x 28 \quad Big white empty looking - with lots of space \quad \uncertain{Bang} barn interesting. Contrast with tall burnt grass - all very blond. \uncertain{Mid} green hill back of barn to R. \quad Robt. W. Huntington, Hartford. Oct. 31. \quad 600 - 1/3 & 400 & Jan. 19, 32. \\
 & & Lombard's House \quad 20 x 28 \quad Blond + very clear. White house, picket fence. Little tree in front of house, \uncertain{trunks} bent to side. Tarred road. \quad \uncertain{in} El Palacio (750) \quad Returns by Whitney in exchange - \quad Resold \ill{} Hopper Retrospective 1950 * \quad Ladies Home Journal. & 750 - 1/3 = 500 . & June 16 1953 \\
 & & Whitney Museums out of First Biennial Exhibition of Contemp. Amer. Sculpture, Water colors + Prints \quad '' Mar 1950 '' Exchanged for Mexican \uncertain{Constructors} of 1946 - 750 - 1/3 Rehn Gallery = 500 & & Feb. 1'' 1934. \\
 & & Dead Tree + Side of Lombard House \quad 20 X 28 \quad Dead tree dark \quad Clifford Frondel - Prof. at Harvard - 750 - 1/3 . \quad 38 \uncertain{Fannald} Drive. \quad Cambridge, Mass. & & June 22, 1950. \\
 & & Roofs of the Cobb Barn \quad 20 X 28 \quad Dark red in shadow. \ill{} barn structures. Bright light on tarred roofs at various angles. One of his best this summer - Interesting for plans. \quad Sold. Encyclopedia Britanica & Rec'd 450 & Ap. 11'', 1944 \\
\end{ledgertable}

\marginal{See note Book III \quad --- in pencil beside the Whitney exchange
row}

\marginal{\uncertain{Cos.} of the Encyclopedia \\ '' Britanica '' - \\ 750. -
10\% to '' \\ - 33 1/3 Gal. \\ p. 78. \quad --- at the right of the Roofs of the
Cobb Barn row}

\hand{jo}{1931 Continued p. 78.}

\note{the 750 in « (750) » of the Lombard's House row is ringed in pencil. The
word between « Resold » and « Hopper Retrospective » will not settle and is
given as \ill{} rather than guessed; « out of » will support it. The \ill{} in
the Roofs of the Cobb Barn description is a word finished by an ink blot, not
paper. « Britanica » is as written, twice. The buyer's street on the Methodist
Church Tower row is offered as « Seoley » and is not read; « Sheaffer » in the
pencil line above and « Sheafer » in the ink line below are as written, adding
to the two forms leaves 54 and 55 already carry. « Constructors » and
« Fannald » are offered rather than read.}

\section{Leaf 76: Etchings at the Rehn Gallery}

\sheet{17740}{76}

\hand{jo}{Water Colors of 1931 continued p. 78.}

\hand{jo}{\emph{Etchings at the Rehn Gallery.}}

\begin{ledgertable}{Y{0.11}Y{0.05}Y{0.48}Y{0.08}Y{0.13}}{ & & & Rec'd. & When rec'd.}
 & 25 & \struck{East Side Interior} & & \\
Oct. 8, 24 & 22 & Evening Wind - bgt. by Mr. Steph. C. Clark & 14.67 & June 11, 25 \\
Mar. 15, 26 & 22 & Aux Fortifications \quad gift to John Clausen, Xmas 27. & & \\
 & 18 & Amer. Landscape & 12. & Nov. 4, 1927. \\
 & 18 & Cat Boat & & \\
 & 18 & The Lonely House & & \\
 & 18 & House by a River \quad Taken for Mod. Mus. Oct. 1933. & & \\
 & 18 & Monhegan Boat - St. Bot. Cl. - 2 Copies & 2 \quad 25.20 & June 15, 26 \\
 & 18 & Cow + Rocks & & \\
 & 25 & Night Shadows & 16.66 & Nov. 4, 1927 \\
 & 18 & The R.R. \quad Farjeon 18 - 1/3 & 12. & Feb. 7, 1929. \\
 & 18 & Girl on a Bridge \quad Farjeon 15 - 1/3 & 10 & Feb. 7, 1929. \\
 & 25 & Night in the Park & 16.66 & Nov. 4, 1927. \\
 & 25 & East Side Interior - St. Bot. Cl. - Ap. 1926. & 17.50 & June 15, 26 \\
 & 15 & House Tops \quad O'Connell \quad 15 - 1/3 & 10. & Feb. 7, 29. \\
 & 18 & Locomotive - sold at St. Bot. Cl. Ap. 1926 & 12.60 & June 15, 26. \\
 & 25 & Evening Wind - sold at St. Bot. Cl. \quad '' & 17.50 & June 15, 26. \\
 & \struck{18} & Night in the L Train \quad Farjeon 18 - 1/3 & 12. & Feb. 7, 1929. \\
 & 18 & R.R. Crossing - drypoint \quad Whitney Retrospective \quad 30 - 1/3 = 20 & & June 22, 1950 \\
 & 18 & Trains + Bathers & 12 & June 5, 29. \\
 & 18 & Maine Coast \quad 25 - 1/3 & 16.66 & June 5, 29. \\
 & 18 & \struck{Bull Fight} \quad 18 - 1/3 & 12. & June 5, 29. \\
 & 18 & Les Deux Pigeons - St. Bot. Cl. & 12.60 & June 15, 26. \\
Nov. \uncertain{24}, 26 & 25 & Evening Wind & 16.66 & May 18, 27 \\
 & \struck{25} & Girl on a Bridge \quad - Taken for Mod. Mus. Oct. 1933. & & \\
 & 25 & Night in the Park \quad Farjeon - 25 - 1/3 & 16.66 & Feb. 7, 29. \\
 & 25 & East Side Interior & 16.66 & Nov. 4, 27. \\
Feb. 3, 27 & & Aux Fortifications \quad for Carnegie group Dec. 1948 - Paid. & & \\
 & & East Side Interior \quad Farjeon 25 - 1/3 & 16.66 & Feb. 7, 1929. \\
 & & The Locomotive \quad O'Connell? \quad 20 - 1/3 & 13.33 & Feb. 7, 1929. \\
 & & Cat Boat \quad Farjeon - 18 - 1/3 & 12. & Feb. 7, 29. \\
 & & R.R. \quad Prof. Saunders - 20 - 1/3 & 13.33 & Feb. 7, 29. \\
Feb. 8, 27 & & Amer. Landscape \quad Farjeon \quad framed? \quad 18 - 1/3 & 12. & Feb. 7, 1929. \\
 & & Evening Wind \quad unframed? & 16.66 & Nov. 4, 27. \\
 & & Night Shadows \quad also \$2 pr frame. & 16.66 & May 18, 27. \\
 & & Night in Park \quad 30 - 1/3 & 20. & June 5, 29. \\
 & & Cow + Rocks \quad Farjeon 18 - 1/3 \quad frame? & 12. & Feb. 7, 29. \\
\end{ledgertable}

\marginal{\uncertain{St. Bo. Joseph's} at \$18 \\ Clark show \quad --- in the
left margin beside the Amer. Landscape row}

\marginal{\struck{framed} \quad --- written vertically beside the brace that
groups the five rows from Feb. 3, 27}

\marginal{4 frames paid for Feb. 7, 1929. \quad --- written vertically up the
left margin beside the last two blocks}

\note{the two blocks headed « Feb. 3, 27 » and « Feb. 8, 27 » are each drawn
together by a brace at the left, and the rows beneath the first line of each
carry no date and no price of their own. Several rows of this leaf carry pencil
that has been erased and is not transcribed. The struck « East Side Interior »
at the head stands above the first ruled line and has no date. In the Monhegan
Boat row a brace joins the « 2 » to the figure beside it, one receipt for the
two copies the title line counts. The day of
« Nov. 24, 26 » is written over another figure and will support « 29 » as
readily. The left margin note beside Amer. Landscape will not settle: « St. Bo.
Joseph's » is offered rather than read, and the abbreviation « St. Bot. Cl. »
that this leaf uses four times elsewhere will support its first half as
readily. « Les Deux Pigeons » is plural here, against
« Les Deux Pigeon » on leaf 40. The right-hand edge of this photograph shows
the facing leaf 77 at the fore-edge; its column is transcribed at sheet 18222.}

\section{Leaf 77: Set of etchings belonging to J. Hopper}

\sheet{18222}{77}

\hand{jo}{19 Etchings to Utica Art Soc. \quad Rudolph Tandler - \quad where
are they now?\\
12 \quad Feb. 25th, 1928\\
1 \quad Mar. 19\\
6 \quad Mar. 27}

\hand{jo}{Set of Etchings belonging to J. Hopper. all framed by Carnegie Art
Inst. when shown there 1937 \quad Left at Rehn Gal. as samples. They all
returned for Whitney Retrospective and now in studio - dedicated to J. H. +
\uncertain{her est.}}

\begin{ledgertable}{Y{0.11}Y{0.05}Y{0.48}Y{0.08}Y{0.13}}{ & & & & }
Feb. 8, 27 (con.) & 18 & House by a River & 12 & May 18, 27 \\
 & & House Tops \quad 15 - 1/3 \quad Farjeon & 10 & Feb. 7, 29 \\
Mar. 11, 27 & 18 & Amer. Landscape \quad for Carnegie group Dec. 1948. & & \\
'' & 25 & Night Shadows \quad + frame = 3. \quad Chase & 16 2/3 & Mar. 9, 28 \\
'' & 18 & Locomotive & 12 & Nov. 4, 1927 \\
'' & 18 & '' & 12 & Mar. 9, 1928 \quad Jackson \\
'' 14 & 25 & Night Shadows & 16.66 & Nov. 4, 1927 \\
Ap. 27 & 25 & East Side Interior & 16 2/3 & Mar. 9, 28 - Davis \\
'' & '' & Evening Wind & 16.66 & May 18, 27 \\
June 2, 27 & 25 & East Side Int. \quad Prof. Saunders & 16.66 & Feb. 7, 29 \\
'' & 18 & Cat Boat \quad Price raised to 25, fall 1927 & 16 2/3 & Mar. 9, 28. Doody \\
Nov. 5, 27 & 25 & Night Shadows & 16 2/3 & Mar. 9, 28 \quad Cunningham \\
Dec. 9, 27 & 25 & East Side Int. \quad Mrs. Hammond. & 20 & Mar. 5, 31. \\
'' & 20 & Locomotive \quad 18 - 1/3 & 12 & Feb. 7, 29 - Farjeon \\
'' & 25 & Cat Boat \quad 30 - 1/3 & & Mar. 5, 31. Hammond \\
'' & 25 & Night in the Park \quad Dec. 1948 for Carnegie group & & \\
'' & 25 & Evening Wind & 16 2/3 & Feb. 7, 29 \quad Farjeon \\
Dec. 16, 27 & 27 & Night Shadows framed & 16 2/3 & Feb. 7, 29 \quad Farjeon \\
Feb. 25, 28 - & 25 & Amer. Landscape \quad Harmon - 30 - 1/3 - Harmon & & Ap. 1, 37 \\
 & 25 & Eastside Interior & & \\
 & 25 & Cat Boat \quad 30 - 1/3. \quad \uncertain{Baumgstein} - June 16, 53. & & \\
 & 25 & Night Shadows \quad 30 - 1/3 - 20 \quad Harmon & & June 16, 53 \\
 & 18 & House by a River & & \\
 & 25 & Evening Wind . Prof. Saunders & 16 2/3 & Feb. 7, 1929 \\
 & 20 & Locomotive \quad 30 - 1/3 & 20 & June 5, 29 \\
 & 18 & Cow + Rocks & & \\
 & 15 & Girl on a Bridge & & \\
Mar. 19, 28 & 25 & Amer. Landscape \quad for Prof. Saunders & & \\
Mar. 27, 28 & & East Side Int. & & \\
'' & & Evening Wind & 20. & Lenox \quad Mrs. A. Sterner \quad Jan. 8, 34. \\
'' & & East Side Int. & & \\
'' & & Evening Wind & & \\
'' & & Amer. Landscape & & \\
\end{ledgertable}

\marginal{+ 2 for frames \quad --- above the House by a River row; framed \quad
--- below its date}

\marginal{For Utica, N.Y. - Utica Art Soc. - Rudolph Tandler. \quad --- written
vertically up the left margin beside the last five rows}

\note{the three counted lines of the heading are gathered by a brace at their
left, and a second brace at their right carries « where are they now? » across
all three; both are drawn and neither is transcribed as a character.}

\note{these are the left of the two runs this leaf carries, read down. The
right-hand run below is not a continuation of it: it was written later, in
another ink, and its dates do not pair with the left's. The five rows from Mar.
27, 28 are drawn together by a brace and are the nineteen etchings the heading
sent to Utica. Several rows carry pencil that has been erased and is not
transcribed; the second Locomotive row is one of them, and its title is a ditto
mark under the row above.}

\begin{ledgertable}{Y{0.12}Y{0.05}Y{0.50}Y{0.16}}{ & & & }
Mar. 27, 28 & 25 & East Side Int. \quad for John Taylor Arms. & 16 2/3 \quad Feb. 7, 1929. \\
Nov. 19, 28 & & Locomotive \quad Dec. 1948 for Carnegie \struck{purchase} & \\
'' '' & 30 & Am. Landscape \quad Boston \quad sent to Calif. & 20. \quad C. H. Bentley \quad Dec. 31, 34 \\
Feb. 6, 29 & 30 & Cat Boat . \quad Mrs. S. Tucker - to go to Boston \quad 30 - 1/3 = 20 & June 5, 29. \\
Mar. 18, 29 & 30 & Locomotive & \\
 & 25 & Maine Coast & 16 2/3 \quad June 5, 29. \\
 & 30 & East Side Int. \quad Tiedtke & 20 - Nov. 20, 36. \\
 & 30 & Boat. & \\
 & 30 & Night in Park. & \\
Oct. 16, 31 & 30 & East Side Interior \quad used by Creative Arts & \\
 & & Cat Boat \quad To Mrs. H. C. Bentley . \quad 20 & Jan. 8, 33. \\
Oct. 23, 33 - & & for Mrs. A. Sterner \quad Evening Wind - already sold by her. & Jan. 8, 34. \\
Dec. 7, 33 & & Amer. Landscape - H. \struck{Whitney Mar.} & Jan. 8 \quad 1933. \\
'' '' & & Even. Wind \quad for Mrs. \uncertain{Joham} & Jan. 8, 34 \\
Ap. 1, 37 & & Evening Wind \quad no mat. \quad Solomon. & Dec. 29, 37 \\
Dec. 29, 37 & & Amer. Landscape. & \\
 & & East Side Interior - Johnson - & Dec. 26, 40. \\
Jan. 20, 41. & & Night in the L train & \\
all 30. & & Rail Road \quad Dec. 1948 Carnegie \struck{purchase} \quad Univ. of Nebraska. & \\
 & & East Side Interior - 50 - 1/3 \quad Craig. & Dec. 29, '47 \\
Feb. 16, 45. \struck{Jan.} & \$50 & East Side Interior \quad \$50 \quad Dec. 1948 Carnegie \struck{purchase} & \\
 & \$30 & Night Shadows \quad \$30 \struck{50} - 1/3 . \quad Des Moines & June 16, 53. \\
June 27, 49 & & Collection of 11 for 300 to Carnegie. See p. 46. & \\
May 27, 51 - & & See p. 35. \quad Rehn Gal. cont. & \\
\end{ledgertable}

\marginal{for Kraushaar \quad when? \quad --- in pencil beside the East Side
Int. / Tiedtke row}

\note{this is the right of the two runs, and it rules four columns where the
left rules five. « Tandler » is settled by the vertical note in the left
margin, against the reading « Taudler » the heading alone would allow. « Joham »
is offered rather than read. The left-hand edge of this photograph shows the
facing leaf 76 at the fore-edge; its column is transcribed at sheet 17740.}

\section{Leaf 78: Continued, Summer 1931, South Truro}

\sheet{17206}{78}

\hand{jo}{Cape Cod.}

\hand{jo}{Continued - Summer 1931 \quad South Truro mostly - Wellfleet nearby.
+ Truro}

\begin{ledgertable}{Y{0.10}Y{0.06}Y{0.44}Y{0.09}Y{0.14}}{ & & & & }
Oct. 1931. & & High Road \quad 20 x 28 \quad Tarred hill road - going into N. Truro. before it dips down. Hills beyond. & & \\
 & & Capt. Kelly's House \quad 20 x 25 \quad Dark picture. Houses smaller than this proportionately. Inky black tree back of house. House has cerulean blue shutters. R.R. track runs in front of house \quad Hill in shadow. Fine tree + shrubbery to R. \quad Sandy road in front + bank for R.R. \quad One of best of 1931 catch. & & Returned to Studio \quad Nov. 12, 36 - \quad frame used for other picture. \\
 & & Freight Car at Truro \quad 14 x 20. \quad Car all in shadow - Sand grey with soot. \quad Clifford Frondel (- Harvard Prof.) \quad 400 - 1/3 = 266 2/3 & & June 22, 1950. \\
 & & The Schuman House \quad 14 x 20 \quad Light on little tree in front of picture stands out against shadow side of house. Nice picture of old house. \quad Mr. Solomon - 400 - 1/3 = 266 2/3. & & Feb. 9, 1937. \\
 & & The Lewis Barn \quad 14 x 20 \quad Bright blue pond to left of barn \quad Pale sand L. lower corner. Dune in back of barn. \quad Mrs. Maguire \quad 400 - 1/3 - & 266 2/3 & May 27, 37 \\
 & & R.R. Warning \quad 14 x 20 \quad Very clear day. Road across R.R. tracks sandy road on way to beach. \quad University Club of St. Louis - \ill{} & 400 - 1/3 - 266 2/3 & Jan. 8, 1934 \\
\end{ledgertable}

\sketch{}

\sketch{}

\sketch{}

\sketch{}

\sketch{}

\sketch{R.R.}

\marginal{from p. 75 \quad --- in pencil above « Continued »}

\marginal{little \quad --- above « of old house » in the Schuman House block,
with a caret}

\marginal{33 \quad --- at the foot of the leaf, below the last line and cut by
the leaf's edge}

\note{each of the six works carries one of Edward Hopper's ink record drawings,
ruled into a box beside the entry; only the last is lettered, with an arrow and
« R.R. » beside the crossing sign, and the rest take an empty argument. The
\ill{} after « University Club of St. Louis » is a faint pencil date, in part
under the sketch box, which will not settle: it will support « Mar. 1930 » and
is not offered as a reading. « Maguire » here is the spelling leaf 73's second
row gives, against « McGuire » on that leaf's last row. A line of erased pencil
stands above the Frondel line of the Freight Car block and is not transcribed.}

\section{Leaf 79: 1931 Cape Cod, and the summers}

\sheet{17175}{79}

\hand{jo}{1931 \quad Cape Cod.}

\begin{ledgertable}{Y{0.10}Y{0.06}Y{0.44}Y{0.09}Y{0.14}}{ & & & & }
1931 & & Wellfleet Road \quad 14 x 20 \quad Tarred road down to the Chiquesset Hotel along beach. A grey day. shaped structure - shacks on R. of road (spread out toward base. \quad Mr. + Mrs. Leslie Gould (Ladies Home Journal) & 400 - 1/3 = 266 2/3 & June 27, 1950 \\
1931 & & Scrub Pines \quad 14 x 20 \quad Dark under tangled droopy pine boughs. White streaks where light strikes rows of pine trunks. Hen coop under boughs middle distance. Pink floor pine needles where sun let in. \quad Fred Palmer bought 3 water colors for 1350. with new frames. & 450 - 1/3 = 300. & Ap. 7 \quad 1956. \\
1931 & & House on Dune Edge \quad 14 x 20 \quad Bright blue sky. Light stucco house with tower effect sitting on top of dune near station at N. Truro. Tower very round, light reflected on underside of round roof. Bushes (bay berry?) in front of house very interestingly massed. \quad Fred Palmer bgt. 3 for 1350 with new frames - & 450 - 1/3 = 300. & Ap. 7 \quad 1956. \\
\end{ledgertable}

\sketch{}

\sketch{}

\sketch{}

\note{a small drawn trapezoid stands on the leaf between « A grey day. » and
« shaped structure » in the Wellfleet Road description, and is the shape the
sentence names. It is a drawing rather than a word and nothing is put in its
place above. Each of the three works carries one of Edward Hopper's ink record
drawings, none of them lettered. « Chiquesset » is as written.}

\hand{jo}{Summers - \qquad See other Ledgers for other water colors of later
date}

\hand{jo}{1923 \quad Gloucester\\
1924\\
1925. New Mexico\\
1926 - Rockland + Cape Anne\\
1927. Two Lights, Cape Elizabeth\\
1928 - Cape Anne\\
1929 \quad Charleston\\
\qquad Cape Elizabeth - (Two Lights)\\
1930 - Cape Cod.\\
1930 - 194\struck{5}2 Cape Cod + trips to Vt. etc.\\
1941 - Trip to Pacific - then back to Cape Cod\\
1942 - Cape Cod -\\
1943 Mexico. no gas to get to Cape - R.R. to Mexico\\
1944 ? gas? Truro}

\marginal{II \quad --- above « Ledgers » in the heading of this list}

\marginal{Vermont 1936 \\ 1937 \\ 1938 \quad --- at the right of the list,
beside the 1930 and 1930--1942 lines}

\note{the list heads no column and is transcribed as the lines stand. A brace
at the left gathers 1923 and 1924 to the single « Gloucester », which stands
beside the first of them and belongs to both. « Cape Anne » is as written,
twice, confirming leaves 56, 57 and 32.}

\section{Leaf 80: Oils, from p. 59}

\sheet{17529}{80}

\hand{jo}{from p. 59.}

\hand{jo}{\emph{Oils}}

\hand{jo}{60'' x 65''\\
Hotel Room - Tall girl, brown hair, reading time table (yellow). Night out
doors, light from overhead. Green carpet, green fat arm chairs, yellow window
shade. Brown oak (immitation mahogany bed + dresser. White spread + pillow
lavender blanket on foot of bed. Black suit case + brown bag. Black hat +
slippers. Very white wall except in shadow at. Extreme left - dull blue grey,
surbase darker;}

\hand{jo}{Ap. 25, 31 \quad Painted in N.Y. Studio \quad In Studio now}

\hand{jo}{Returned from \uncertain{Cohn} Gallery. June 30, 49. \quad Sold by
Rehn Gallery to Mr. + Mrs. N. B. Spingold \quad 5500 - 1/3 = 3666 2/3 \quad
Dec. 27, 1955.}

\sketch{}

\marginal{in pink shirt \quad --- above « reading » in the Hotel Room
description}

\marginal{12 E. 77 \quad --- in pencil below « Mr. + Mrs. »}

\marginal{She \ill{} \uncertain{Frances} - dress shop \quad --- in pencil above
the Spingold line, with « Com » written across it in ink}

\hand{jo}{The Camel's Hump\\
32 X 50\\
Oct. 1, 1931.\\
S. Truro.}

\hand{jo}{S. Truro hills back of the church \quad variety of greens. Bare
saddle shaped dune (Indian Camp ground) on sky line. Blue sky white clouds.
Late after noon - Hills in back in shadow, all of them covered with foliage
except Camel Backs. that bare. Foreground boggy with lots of green + blue green
foliage.}

\hand{jo}{(Arthur N.)\\
Mrs. Eleanor Pack \quad Princeton \quad Fall 1931\\
\struck{The} Purchased from Pack by Edw. W. Root - Feb. 1936.}

\begin{ledgertable}{Y{0.20}Y{0.14}Y{0.14}Y{0.16}Y{0.18}}{ & & & & }
2500 + returned water color - Shore Acres = 2500. & & & & \\
1'' payment & 500. & - 1/3 = & 333 1/3 & Jan. 19, 32. \\
2'' '' & 1700 & - 1/3 = & 1133 1/3 & Dec. 29, 32 or Jan. 33 \\
3'' '' & 300 & - 1/3 = & 200 & Jan. 5, 1934. \\
\end{ledgertable}

\sketch{}

\hand{jo}{N.Y., New Haven + Hartford\\
32 X 50\\
Oct. 1, 1931.\\
See opposite page.\\
S. Truro.}

\hand{jo}{Green hill in deep shadow with streaks of sunlight across top - +
outlining edges of some tree branches - trees in shadow. Sun in back. Foreground
sand in shadow at near side of track. sky not very bright blue. White clouds -
greyish ones.}

\hand{jo}{Indianapolis Mus. After Nov. 1932. paid R. \quad 2500 \quad John
Herron Art Institute, Indianapolis \quad 2500 - 1/3 = 1666 2/3 \quad June 1,
32.}

\sketch{}

\hand{jo}{Room in Brooklyn\\
29 x 34\\
Feb. 2, 1932.\\
Painted in N.Y. Studio}

\hand{jo}{Green carpet with bright yellow patch of sun light from right window.
Red brick houses opposite. White vase. blue green cloth under it. Pinky
flowers, lots of leaves. Green window shades. Dark red table cover to L. Woman
black hair - blue grey dress. Raw sienna wood work + rocking chair. Blue sky.
(Window curtains down to be washed)}

\hand{jo}{Boston Museum of Fine Arts \quad 1800 - 1/3 = 1200 \qquad May 14,
1935}

\marginal{Arts \quad --- above « Institute » in the John Herron line}

\marginal{sky. \quad --- above « Blue » in the last line of the Room in Brooklyn
description}

\note{this leaf and leaf 81 rule nothing and head nothing: each work stands as a
block with its title, size and date at the left, one of Edward Hopper's large
ink record drawings across the middle, the description in a column to the right
of it, and the sale written across the foot. The Camel's Hump payments are the
one place the leaf sets figures in columns and are given as a table. The « 2500
+ returned water color » is the exchange the Shore Acres block of leaf 71
records from the other side. « immitation » is as written. « Cohn Gallery » is
offered rather than read; the same name stands again on leaf 81 and is no
clearer there. The pencil above the Spingold line will not settle beyond
« dress shop », and the word offered as « Frances » may be « frames ».}

\section{Leaf 81: Barber Shop and Room in N.Y.}

\sheet{18310}{81}

\hand{jo}{Barber Shop\\
5 x 6 1/2 ft.\\
Dec. '31\\
Painted in N.Y. Studio}

\hand{jo}{Brass rail, white walls, blue green table top. Manicurer girl in
light green dress, white collar reading the Graphic. Brown hair, flesh colored
stockings. Barber \struck{two} reddish brown trousers, white coat. Old goldish
colored curtain at window. Barber pole shows outside window. + shadow side of
buildings opposite above sidewalk +}

\hand{jo}{Returned from \uncertain{Cohn} Gallery, June 30 \struck{1948} 1949 .
Back to Rehn Gallery, after treatment by \ill{}, June 1954. \quad Bgt. by Roy
Neuberger \quad \$4000 - 1/3 = \$2666 2/3 \quad Sept. 2, 1954.}

\sketch{}

\marginal{120 Bdwy, N.Y. \quad --- in pencil above « Bgt. by Roy Neuberger »;
Mr. + Mrs. \quad --- in pencil before it}

\hand{jo}{Room in N.Y.\\
29'' X 36''\\
Feb. \uncertain{24} '32\\
Painted in N.Y. Studio}

\hand{jo}{Outside night. Inside bright green walls, oak wood work door, oak
table. Girl in bright red, sitting front, head + shoulders twisted side ways to
face piano, picking at keys with one finger. Arms + neck bare, very white
flesh, dark hair, profile in shadow but light on cheek + neck. Blond man in
shirt sleeves reading newspaper. Pink upholstered chair. Dark red lamp shade
showing sticking out from piano at top of girl's head. 4 blocks of masonry
outside at left + pillar \ill{}\\
\ill{}ish reflected light. Under window black.}

\hand{jo}{University of Nebraska \quad at Lincoln, Neb. \quad (Exhibit
\struck{run} by Walker Gallery.) \quad 1200 - 1/3 = 800. \quad May
\uncertain{20}, 36}

\sketch{}

\marginal{managed by \quad --- above « run by » in the Nebraska line}

\clipping{a halftone of the Camel's Hump oil --- bare dunes under a wide sky,
the saddle-shaped hill on the sky line --- printed on a cream leaf and titled
« Camel's Hump » in bold above it, pasted flat over the lower two thirds of the
leaf.}

\hand{jo}{Nov. 2, \ill{}\\
Justi\ill{}\\
the Sk\ill{}\\
Nov. 13 -\\
* French\ill{}\\
the Cam\ill{}}

\note{these six pencil lines stand at the left of the clipping with an arrow
drawn up from them, and every one of them is cut off where the clipping begins.
Each \ill{} there means « under paper », and unlike leaves 66, 71 and 75 there
is no second photograph of this leaf anywhere in the Whitney's sequence: the
clipping is pasted flat, not hinged, and what it covers is not recoverable from
these photographs. The word broken across the last two lines of the Room in N.Y.
description will not settle in either half. A brown stain lies over the
conservator's name in the Barber Shop line and only its last letters show; the
name is not offered. The day of « Feb. 24 '32 » is written over another figure,
and the day of « May 20, 36 » is cut by the fore-edge.}

\keywords{medium:watercolours, medium:oils, medium:etchings, place:South Truro,
place:North Truro, place:Wellfleet, place:Provincetown, place:Cape Cod,
place:Gloucester, place:Rockland, place:Essex, place:Cape Ann,
dealer:Frank K. M. Rehn, dealer:E. Weyhe, dealer:Kraushaar,
person:Robert Wheelwright, person:Edward W. Root, person:Clifford Frondel,
person:Robert W. Huntington, person:Fred Palmer, person:Roy Neuberger,
person:N. B. Spingold, person:Eleanor Pack, person:Alfred L. Hydeman,
person:Leslie Gould, person:Mrs. Albert Sterner, person:Farjeon,
person:A. Percy Saunders, person:John Taylor Arms,
person:David Rockefeller, person:Mrs. Bliss Parkinson,
collection:Wadsworth Atheneum, collection:Fogg Art Museum,
collection:Boston Museum of Fine Arts,
collection:John Herron Art Institute, collection:Whitney Museum,
collection:Carnegie Institute, collection:University of Nebraska,
society:Utica Art Society, society:St. Botolph's Club,
feature:dealer commissions, feature:exchanges, feature:resales,
feature:record sketches, feature:pasted clippings, feature:unlisted works}

\end{document}
